\documentclass[acmsmall,screen]{acmart}
\citestyle{acmauthoryear}

%% --- ACM metadata --------------------------------------------------------
\acmJournal{PACMPL}
\acmVolume{1}
\acmNumber{HASKELL}
\acmYear{2026}
\acmDOI{}

%% --- Packages ------------------------------------------------------------
\usepackage{listings}
\usepackage{amsmath}
\let\Bbbk\relax
\usepackage{amssymb}
\usepackage{stmaryrd}   % \llbracket, \rrbracket
\usepackage{mathtools}
\usepackage{booktabs}
\usepackage{microtype}
\usepackage{enumitem}

%% --- Fix headheight for fancyhdr -----------------------------------------
\setlength{\headheight}{18.67603pt}
\addtolength{\topmargin}{-5.67603pt}

%% --- Haskell listings style ----------------------------------------------
\lstdefinelanguage{Haskell}{
  keywords={data, type, class, instance, where, let, in, do, if, then, else,
            case, of, import, module, deriving, infixl, infixr, Not, Or, And,
            Seq, Bot, Epsilon, Single, Star, Atom, Wildcard,
            GuardedRE, WRE, Semiring, Prob, Tropical, PTrue, PFalse},
  keywordstyle=\color{blue!70!black}\bfseries,
  sensitive=true,
  comment=[l]{--},
  morecomment=[s]{\{-}{-\}},
  commentstyle=\color{gray}\itshape,
  stringstyle=\color{orange!80!black},
  morestring=[b]",
  literate=
    {<>}{{$\diamond$}}1
    {/\\}{{$\wedge$}}1
    {\\\\}{{$\setminus$}}1
    {->}{{$\to$}}2
    {=>}{{$\Rightarrow$}}2
    {forall}{{$\forall$}}1
    {>>}{{$>\!\!>$}}2
    {>>=}{{$>\!\!>\!\!=$}}3,
}
\lstset{
  language=Haskell,
  basicstyle=\ttfamily\small,
  columns=flexible,
  keepspaces=true,
  frame=none,
  xleftmargin=1em,
  breaklines=true,
  breakatwhitespace=true,
}

%% --- Math macros ---------------------------------------------------------
\newcommand{\RE}{\mathsf{RE}}
\newcommand{\GRE}{\mathsf{GRE}}
\newcommand{\WRE}{\mathsf{WRE}}
\newcommand{\SL}{\mathsf{SL}}
\newcommand{\Eff}{\mathsf{Effectful}}
\newcommand{\sep}{*}
\newcommand{\wand}{-\!\!*}
\newcommand{\pre}{\mathit{pre}}
\newcommand{\post}{\mathit{post}}
\newcommand{\fut}{\mathit{future}}
\newcommand{\ret}{\mathit{ret}}
\newcommand{\anything}{\mathtt{anything}}
\newcommand{\finally}[1]{\mathtt{finally}(#1)}
\newcommand{\previously}[1]{\mathtt{previously}(#1)}
\newcommand{\shuffle}{\mathbin{\sqcup\!\!\!\sqcup}}
\newcommand{\never}[1]{\mathtt{never}(#1)}
\newcommand{\noUntil}[2]{\mathtt{noUntil}(#1,\,#2)}
\newcommand{\derives}[2]{\partial_{#1}(#2)}
\newcommand{\norm}[1]{\lceil #1 \rceil}
\newcommand{\sem}[1]{\llbracket #1 \rrbracket}
\newcommand{\univ}{\neg\varnothing}

\begin{document}

%% --- Title / Authors -----------------------------------------------------
\title{Effectful Computations with Future Conditions:\texorpdfstring{\\}{}
       A Monadic Approach to Temporal Specification}

\author{Author Name}
\affiliation{%
  \institution{Institution}
  \city{City}
  \country{Country}
}
\email{author@institution.edu}

\begin{abstract}
Pre- and post-conditions are the standard vocabulary for modular program
specification, yet they are silent about obligations that span multiple
calls.
In real Haskell code such obligations are ubiquitous: a successful TLS
handshake must eventually be followed by \lstinline{bye} before the
context is discarded; a taken \lstinline{MVar} must eventually be put
back on \emph{every} code path; a \lstinline{beginTx} must eventually reach
\lstinline{commit} or \lstinline{rollback}; a submitted job must complete
with probability at least~$p$.
No single pre/post pair can express these properties on its own, because the
balancing operation lies in an unknown future context.

Building directly on \citet{song2025futurecond}'s proposal of future
conditions as a first-class extension to pre/post contracts, we present
\emph{Pledge}, a Haskell library that realises this idea as a monad:
annotations that propagate obligations forward until they are discharged.
The \lstinline{Pledge} monad carries a data-dependent field
$\mathtt{future} :: a \to \mathit{eff}$, allowing obligations to name the
specific resource handle an operation returns.
Monadic bind partially discharges the first computation's future
condition against the continuation's postcondition via
\emph{subtraction}; the residual is conjoined with the continuation's
own future condition.
Pledge is a \emph{specification} library: it computes and reports residual
obligations on pure annotation values, it does not instrument or guard a
running \lstinline{IO} computation, and the accompanying monad-law proof
covers the $(\ret,\post,\fut)$ triple, not the \lstinline{pre} field, which
is carried as a separate contract-checking annotation.

The design is parameterised over a \lstinline{Composable} algebra with
concatenation~($\diamond$), conjunction~($\wedge$), subtraction~($\setminus$),
and identities \lstinline{empty} and \lstinline{universe}.
We present four instances: trace protocols via extended regular
expressions~($\RE$, with Brzozowski-derivative subtraction and a direct
LTL$_f$ embedding); Presburger-guarded traces~($\GRE$, discharged by
Z3); semiring-weighted obligations~($\WRE$, covering probabilistic and
min-cost scenarios); and separation-logic heap predicates~($\SL$, with
magic wand as subtraction).
All four share a single bind formula, demonstrating that trace-based,
arithmetic, quantitative, and heap-ownership obligations are uniformly
expressible within one algebraic interface, and we give a soundness
theorem for the $\RE$ instance connecting a fully-discharged residual to
satisfaction of the obligation by every terminating trace.
\end{abstract}

\keywords{temporal specification, monadic effects, regular expressions,
          complement, Brzozowski derivatives, Presburger arithmetic,
          semiring weights, separation logic, resource lifecycle, Haskell}

\begin{CCSXML}
<ccs2012>
  <concept>
    <concept_id>10011007.10011006.10011008.10011009.10011012</concept_id>
    <concept_desc>Software and its engineering~Functional languages</concept_desc>
    <concept_significance>500</concept_significance>
  </concept>
  <concept>
    <concept_id>10011007.10011006.10011008.10011024.10011028</concept_id>
    <concept_desc>Software and its engineering~Data types and structures</concept_desc>
    <concept_significance>300</concept_significance>
  </concept>
  <concept>
    <concept_id>10011007.10011006.10011041</concept_id>
    <concept_desc>Software and its engineering~Software verification and validation</concept_desc>
    <concept_significance>300</concept_significance>
  </concept>
</ccs2012>
\end{CCSXML}
\ccsdesc[500]{Software and its engineering~Functional languages}
\ccsdesc[300]{Software and its engineering~Data types and structures}
\ccsdesc[300]{Software and its engineering~Software verification and validation}

\maketitle

%% =========================================================================
\section{Introduction}
\label{sec:intro}
%% =========================================================================

Pre- and post-conditions are the workhorse of modular program specification.
A pre-condition constrains the state in which an operation may be invoked;
a post-condition describes what holds immediately after it returns.
Yet a great many correctness properties are not local to a single call
boundary: they \emph{span} a sequence of operations, with arbitrary
computation in between.

Consider three scenarios from the Haskell ecosystem:
i) \emph{TLS lifecycle.}
A successful \lstinline{handshake} must eventually be followed by \lstinline{bye}
before the context is discarded.
A post-condition on \lstinline{handshake} records success, but cannot enforce
that \lstinline{bye} will be called, because the discharge point lies in an
unknown future context.
The \lstinline{bracket} combinator sidesteps this by lexically scoping
the context, but cannot express richer message-exchange obligations or
handle contexts passed across function boundaries.
ii) \emph{MVar discipline.}
The \lstinline{takeMVar} operation removes a value, leaving the variable
empty; the caller must restore it with \lstinline{putMVar} on every exit path.
A post-condition records the empty state but cannot require that a future
\lstinline{putMVar} on the \emph{same} variable will occur.
iii) \emph{Database savepoints.}
The \lstinline{transactionSave} function opens a fresh transaction that
must itself be committed or rolled back; pre/post-conditions on the call
cannot carry this obligation forward.

\paragraph{Existing remedies and their limits.}
The \lstinline{bracket} function lexically scopes cleanup but cannot
propagate obligations beyond the lexical scope.
Both \lstinline{ResourceT}~\cite{snoyman2011resourcet} and
\lstinline{Acquire}/\lstinline{Managed} thread cleanup actions through
every bind but model only the cleanup half of an obligation, discharging
it uniformly at scope exit rather than by application-specific operations.
Linear Haskell~\cite{bernardy2018linear} enforces ``consume exactly
once'' but cannot express branching discharge obligations or quantitative
commitments.
Parameterised and indexed monads~\cite{atkey2009} encode protocol states
at the type level but require \lstinline{RebindableSyntax} or
\lstinline{QualifiedDo}, produce opaque type errors, and are confined to
a single domain per design.

\paragraph{Relation to prior future-condition proposals.}
The idea of a future condition itself is not new here: \citet{song2025futurecond}
introduce future conditions as a first-class extension to pre/post-condition
contracts, formalising the notion that an operation may discharge temporal
obligations onto its continuation, and \citet{song2024provenfix} apply
regular-expression-encoded temporal properties to guide automated program
repair in ProveNFix. This paper's contribution is not the concept of a
future condition but a \emph{library realisation} of it: a concrete
\lstinline{Pledge} monad, a derivative-based subtraction mechanism that
implements the discharge relation algorithmically rather than declaratively,
and four \lstinline{Composable} instances --- traces, Presburger-guarded
traces, semiring-weighted traces, and separation-logic heaps --- that the
formal proposal of \citet{song2025futurecond} leaves uninstantiated. We
return to a detailed comparison in \S\ref{sec:related}.

\paragraph{Specification, not enforcement.}
Before going further, we state plainly what \lstinline{Pledge} is and is not,
since this distinction matters for every example that follows.
\lstinline{Pledge} is a \emph{specification} library: a \lstinline{Pledge}
value is a pure Haskell record of annotations (\lstinline{pre},
\lstinline{post}, \lstinline{future}) that describes what a computation
should do. Composing \lstinline{Pledge} values via \lstinline{>>=} computes
a residual annotation; a programmer or a separate tool inspects that residual
\emph{before, or independently of,} running any real effect. \lstinline{Pledge}
does not instrument \lstinline{IO}, does not intercept asynchronous
exceptions, and does not by itself guarantee that a real running program
obeys the obligations its specification names --- it tells you, given a
specification term, whether the obligations it encodes are met by every
trace consistent with that term. \S\ref{sec:effects} discusses concretely
how a \lstinline{Pledge} specification is intended to relate to a real
effectful computation.

\paragraph{Intended use.}
Across our examples, a \lstinline{Pledge}-annotated primitive
(\lstinline{tlsHandshake}, \lstinline{takeMVar}, \lstinline{malloc}) is
written once, by the author of a library or an API, exactly as a Hoare-style
pre/post-condition would be. The \emph{consumer} of a library is the
programmer composing these primitives with \lstinline{>>=}; that
programmer's obligation is only to inspect the residual \lstinline{pre} and
\lstinline{future} of the composed term, not to write temporal
specifications themselves. This division of labour matches ordinary
contract-based programming: application programmers consume contracts,
library authors write them. The four \lstinline{Composable} instances play
different roles in this picture. $\RE$ is the primary, end-to-end usable
instance: subtraction is fully computed by derivatives, and a residual of
$\univ$ is a decidable certificate that every trace-based obligation is met
(Theorem~\ref{thm:soundness}). $\GRE$ is equally usable, discharging the
arithmetic component via an SMT call. $\WRE$ is usable for the quantitative
obligations it targets (probability, cost), though, as
\S\ref{sec:wre-instance} discusses, it deliberately gives up the
complement operator and hence the LTL$_f$ embedding. $\SL$, by contrast, is
included primarily as evidence that the \lstinline{Composable} interface
generalises beyond trace-shaped obligations to heap ownership: as
\S\ref{sec:sl-instance} makes explicit, its \lstinline{subtraction} records
the magic wand \emph{symbolically} and does not invoke a solver, so a
non-\lstinline{universe} SL residual identifies a candidate obligation for
a programmer or an external tool (e.g.\ Smallfoot or Z3) to check, rather
than a machine-checked verdict on its own.

\paragraph{Our approach.}
We propose the \lstinline{Pledge} monad as a lightweight, compositional
mechanism for specifying such obligations, requiring no language
extensions beyond standard Haskell type classes.
A value of type \lstinline{Pledge eff a} carries three annotations:
a precondition $\mathtt{pre} :: \mathit{eff}$,
a postcondition $\mathtt{post} :: \mathit{eff}$,
and a data-dependent future condition
$\mathtt{future} :: a \to \mathit{eff}$.
Because \lstinline{future} ranges over the return value, an obligation
can name the actual resource handle produced at runtime: \lstinline{tlsHandshake}
can declare that \emph{whichever} \lstinline{Context} it returns must
eventually receive \lstinline{bye}.
Monadic bind partially discharges the first computation's future
via \lstinline{subtraction} against the continuation's postcondition;
the residual is conjoined with the continuation's own future condition.
The user annotates individual operations only; the monad propagates
obligations automatically.

The \lstinline{Composable} typeclass abstracts the specification algebra
into five operations: concatenation~($\diamond$), conjunction~($\wedge$),
subtraction~($\setminus$), and their respective identities
\lstinline{empty} and \lstinline{universe}. It admits four independent
instances:

\begin{itemize}[wide, labelindent=10pt]
  \item \emph{(Trace protocol, $\RE$)}
        Extended regular expressions over events, with complement as a
        first-class constructor. Subtraction is implemented via
        Brzozowski derivatives; LTL$_f$ operators embed directly as RE
        terms without automaton construction. This instance handles
        \lstinline{takeMVar}/\lstinline{putMVar} discipline, TLS
        handshake/bye pairing, transaction lifecycle, and lock protocols.

  \item \emph{(Arithmetic bounds, $\GRE$)}
        Each RE is paired with a Presburger arithmetic predicate over
        program variables, checked by an SMT solver (Z3 via SBV). A single annotation simultaneously
        enforces a protocol ordering \emph{and} a numeric bound,
        catching, for example, counter overflow before any execution
        takes place.

  \item \emph{(Quantitative obligation, $\WRE$)}
        Boolean membership is generalised to a semiring weight.
        The probability semiring captures obligations of the form
        ``this job completes with probability at least~$p$'';
        the tropical semiring gives minimum-cost scheduling obligations.
        A weight of \lstinline{szero} in the residual signals an
        unfulfilled quantitative commitment.

  \item \emph{(Heap ownership, $\SL$)}
        Separation-logic heap predicates, with separating conjunction as
        sequential composition and magic wand as subtraction. Future
        conditions describe what heap state must hold when the
        computation completes; a leaked \lstinline{Cell} in the residual
        post names the unreleased ownership precisely.
\end{itemize}

\paragraph{Central insight.}
The unifying mechanism is \emph{Brzozowski-derivative subtraction as
obligation propagation}.
When a computation $e$ is sequenced with a continuation $f$, the
continuation's postcondition $\post(e_2)$ is subtracted from $e$'s
pending future obligation $\fut(e)$ using the derivative:
for each event $a$ that $\post(e_2)$ can begin with,
$\partial_a(\post(e_2)) \setminus \partial_a(\fut(e))$ advances both
sides in lock-step, consuming the obligation event by event.
The residual is what remains unmet.
The insight is that this derivative-based subtraction, originally
introduced for regular-language parsing, doubles as a compositional
discharge mechanism for temporal obligations in a monadic setting.
The four \lstinline{Composable} instances are four realisations of the
same pattern: RE subtraction via derivatives; $\GRE$ subtraction via
derivatives with a Presburger conjunction; $\WRE$ subtraction with
semiring weights; and SL subtraction via the magic wand.
All four share a single bind formula, and the diagnostic residual is
always the same: anything not equal to \lstinline{universe} names an
unmet obligation precisely, and for the $\RE$ instance this diagnostic
is backed by the soundness theorem of \S\ref{sec:soundness}.

\paragraph{Contributions.}
\begin{enumerate}
  \item We define the \lstinline{Pledge} monad (\S\ref{sec:effectful}),
        parameterised over a \lstinline{Composable} algebra, with
        $\mathtt{pre} :: \mathit{eff}$,
        $\mathtt{post} :: \mathit{eff}$, and a data-dependent
        $\mathtt{future} :: a \to \mathit{eff}$ field, explain how this pure
        specification layer relates to real effectful code
        (\S\ref{sec:effects}), and verify that
        the monad laws hold for the $(\ret, \post, \fut)$ triple
        (\S\ref{sec:monad-laws}).

  \item We instantiate the algebra to extended regular expressions over
        events with complement, using Brzozowski derivatives to implement
        subtraction (\S\ref{sec:re}), show that LTL$_f$ operators
        embed directly as RE terms (\S\ref{sec:ltl}), and prove a
        soundness theorem connecting a discharged residual to
        satisfaction of the obligation on every terminating trace
        (\S\ref{sec:soundness}).

  \item We extend the RE instance to a guarded form ($\GRE$,
        \S\ref{sec:guardedre-instance}), pairing each RE with a
        Presburger arithmetic predicate over program variables,
        evaluated by an SMT solver (Z3 via SBV) at membership-check
        time.

  \item We provide a weighted RE instance ($\WRE$,
        \S\ref{sec:wre-instance}), parameterised over a semiring, that
        generalises Boolean membership to a quantitative weight,
        recovering probabilistic and min-cost future conditions as
        special cases.

  \item We provide a fourth, independent instantiation to
        separation-logic heap predicates (\S\ref{sec:sl-instance}),
        demonstrating that the \lstinline{Composable} interface extends
        beyond trace-based specifications, while being explicit that this
        instance records obligations symbolically rather than discharging
        them against a solver.

  \item We demonstrate the approach on use cases across all four
        instances (\S\ref{sec:examples}), showing that resource leaks,
        ordering violations, arithmetic overflows, probabilistic
        obligations, and heap-ownership errors all surface as
        non-trivial residual annotations, favouring examples with
        non-trivial future conditions throughout.
\end{enumerate}

%% =========================================================================
\section{Background}
\label{sec:background}
%% =========================================================================

\subsection{Extended Regular Expressions over Events}
\label{sec:re}

Events are the atomic units of observation; each carries a name and a list of typed arguments.

\begin{lstlisting}
data Term  = Str String | Num Int | List [Term]
data Event = Atom String [Term]   -- concrete event, e.g. Atom "send" [Num 1]
           | Wildcard             -- pattern matching any single event
data RE = Bot | Epsilon | Single Event
        | Seq RE RE | Or RE RE | And RE RE | Star RE | Not RE
\end{lstlisting}

The specification language is an extended regular expression type
that adds complement as a first-class constructor:
\[
  r \;::=\; \varnothing \mid \varepsilon \mid e \mid r_1 \cdot r_2
          \mid r_1 \vee r_2 \mid r_1 \wedge r_2 \mid r^* \mid \neg r
\]
The constructors are standard.
The empty language $\varnothing$ (\lstinline{Bot}) accepts no trace.
The empty-trace language $\varepsilon$ (\lstinline{Epsilon}) accepts only
the empty sequence.
A single-event term $e$ (\lstinline{Single Event}) matches exactly one
event; \lstinline{Wildcard} is a pattern that matches any single event.
Concatenation $r_1 \cdot r_2$ (\lstinline{Seq}) accepts a trace if it
splits into a prefix in $L(r_1)$ followed by a suffix in $L(r_2)$.
Union $r_1 \vee r_2$ (\lstinline{Or}) accepts traces in either language.
The Kleene star $r^*$ (\lstinline{Star}) accepts zero or more repetitions
of traces from $L(r)$; $\Sigma^*$ is the set of all finite event sequences
over the alphabet $\Sigma$ (the universal language).
Intersection $r_1 \wedge r_2$ (\lstinline{And}) could be derived as
$\neg(\neg r_1\vee\neg r_2)$ via De Morgan; we instead include \lstinline{And}
as a primitive constructor (matching the code in
\S\ref{sec:re}) so that its derivative can be computed directly, without
paying for three extra \lstinline{Not} nodes on every conjunction.
Complement $\neg r$ (\lstinline{Not}) is the set of all traces in $\Sigma^*$
\emph{not} in $L(r)$.
Moreover, we use the following five derived forms throughout:

\begin{align*}
  \anything &\;\triangleq\; \neg\varnothing
            &&\text{(the universal language } \Sigma^*\text{)} \\
  \finally{e} &\;\triangleq\; \anything \cdot e \cdot \anything
            &&\text{($e$ must occur eventually)} \\
  \mathtt{globally}(e) &\;\triangleq\; e^*
            &&\text{($e$ at every step)} \\
  \mathtt{never}(e) &\;\triangleq\; \neg\finally{e}
            &&\text{($e$ must never occur)} \\
  \mathtt{noUntil}(e,\,g) &\;\triangleq\;
            \neg\bigl((\Sigma\setminus\{g\})^* \cdot e \cdot \anything\bigr)
            &&\text{($e$ must not occur before $g$)}
\end{align*}

\noindent
$\mathtt{noUntil}(e,\,g)$ is nullable (the empty continuation satisfies it)
and its derivative with respect to $g$ is $\Sigma^*$, meaning once $g$ occurs
the restriction on $e$ is lifted.
This is the key combinator for \emph{conditional} safety: $e$ is forbidden
until $g$ re-enables it.

\subsection{Algebraic Complement and Brzozowski Derivatives}
\label{sec:deriv}

The Brzozowski derivative~\cite{brzozowski1964} $\partial_a(r)$
computes the language of continuations after reading event $a$.
The key law for complement is:

\begin{equation}
  \partial_a(\neg r) \;=\; \neg\bigl(\partial_a(r)\bigr)
  \label{eq:deriv-not}
\end{equation}

The nullability function $\nu(r)$ ($= \mathtt{True}$ iff
$\varepsilon \in L(r)$) satisfies $\nu(\neg r) = \neg\,\nu(r)$.
Together, equations~\eqref{eq:deriv-not} and the De Morgan laws

\[
  \neg(r_1\vee r_2) = \neg r_1\wedge\neg r_2
  \qquad
  \neg(r_1\wedge r_2) = \neg r_1\vee\neg r_2
\]

allow complement to be propagated and simplified \emph{purely
algebraically}, without building a DFA.
The normaliser additionally applies $\neg\neg r = r$,
$\neg\varnothing = \univ$, and $\neg\univ = \varnothing$, and the same
normaliser is reused to simplify $\mathtt{Not}$ nodes arising from
$\mathtt{never}$ and $\mathtt{noUntil}$.
We define the nullable and derivative functions in Haskell as follows:

\begin{lstlisting}
nullable :: RE -> Bool
nullable (Not r)     = not (nullable r)      -- nu(Not r) = not(nu(r))
nullable (Star _)    = True
nullable (Seq r1 r2) = nullable r1 && nullable r2
nullable (Or  r1 r2) = nullable r1 || nullable r2
nullable (And r1 r2) = nullable r1 && nullable r2
nullable Epsilon     = True
nullable _           = False

derivative :: Event -> RE -> RE
derivative e (Not r)      = Not (derivative e r)   -- d_a(Not r) = Not(d_a(r))
derivative e (Single p)   = if matchEvent e p then Epsilon else Bot
derivative e (Seq r1 r2)
    | nullable r1 = Or (Seq (derivative e r1) r2) (derivative e r2)
    | otherwise   = Seq (derivative e r1) r2
derivative e (Or  r1 r2)  = Or  (derivative e r1) (derivative e r2)
derivative e (And r1 r2)  = And (derivative e r1) (derivative e r2)
derivative e (Star r)     = Seq (derivative e r) (Star r)
derivative _ _            = Bot
\end{lstlisting}

The \lstinline{firstWith} function returns the set of events (drawn from
a supplied finite alphabet $\Sigma_0$) that can begin a word in $L(r)$,
driving the subtraction algorithm; \lstinline{first} in the prose below
always refers to this same function, applied to the alphabet of the
expression itself (\lstinline{first r = firstWith (atoms r) r}).
The complement case requires particular care.
Naively writing \lstinline{firstWith alph (Not r) = firstWith alph r} is
\emph{unsound}: $L(\neg r)$ can begin with events \emph{not} in
$\mathit{first}(r)$.
The correct characterisation is:

\begin{equation}
  e \in \mathit{first}(\neg r,\,\Sigma_0)
  \;\iff\;
  e \in \Sigma_0
  \;\wedge\;
  \lceil\partial_e(r)\rceil \neq \univ
  \label{eq:first-not}
\end{equation}

That is, $e$ can open a word in $L(\neg r)$ if and only if, after
consuming $e$, the residual $\partial_e(r)$ does \emph{not} cover the
entire language: there exists some continuation that remains outside
$L(r)$.
The condition $\lceil\partial_e(r)\rceil \neq \univ$ checks exactly this.
If the normalised derivative equals the universal language $\Sigma^*$,
then every continuation after $e$ belongs to $L(r)$; consequently, no
continuation after $e$ belongs to $L(\neg r)$, and $e$ cannot open any
word in that language.



\begin{lstlisting}
atoms :: RE -> [Event]
atoms (Single Wildcard) = []
atoms (Single e)        = [e]
atoms (Seq r1 r2)       = atoms r1 `union` atoms r2
atoms (Or  r1 r2)       = atoms r1 `union` atoms r2
atoms (And r1 r2)       = atoms r1 `union` atoms r2
atoms (Star r)          = atoms r
atoms (Not r)           = atoms r
atoms _                 = []

firstWith :: [Event] -> RE -> [Event]
firstWith _    Bot               = []
firstWith _    Epsilon           = []
firstWith _    (Single e)        = [e]
firstWith alph (Seq r1 r2)
    | nullable r1                = firstWith alph r1 `union` firstWith alph r2
    | otherwise                  = firstWith alph r1
firstWith alph (Or  r1 r2)       = firstWith alph r1 `union` firstWith alph r2
firstWith alph (And r1 r2)       = [e | e <- firstWith alph r1,
                                        e `elem` firstWith alph r2]
firstWith alph (Star r)          = firstWith alph r
firstWith alph (Not r)           = [e | e <- alph,
                                        not (isTotal (normalize (derivative e r)))]
  where isTotal (Not Bot) = True; isTotal _ = False
\end{lstlisting}

\vspace{3mm}
The alphabet $\Sigma_0$ is collected from an RE via \lstinline{atoms},
which recursively gathers all concrete (\lstinline{Atom}) events,
excluding \lstinline{Wildcard}.
Inside \lstinline{subtraction} (\S\ref{sec:re-instance}), both operands
contribute to the alphabet, so that events appearing only under a
complement on one side are still enumerated correctly by
Equation~\eqref{eq:first-not}.


\begin{table}[!b]
\caption{Selected \lstinline{firstWith} evaluations with \lstinline{Not}.
         $\Sigma_0$ is the explicitly supplied alphabet.
         Rows marked $\dagger$ produce a \emph{strict} subset of $\Sigma_0$.}
\label{tab:firstwith-not}
\begin{tabular}{llll}
\toprule
RE $r$ & $\Sigma_0$ & \lstinline{firstWith} $\Sigma_0\;(\neg r)$ & Rationale \\
\midrule
$a$                     & $\{a,b\}$   & $\{a,b\}$  & $\partial_a(a)=\varepsilon$, $\partial_b(a)=\varnothing$; neither $\Sigma^*$ \\
$a \cdot b$             & $\{a,b,c\}$ & $\{a,b,c\}$ & all derivatives non-$\Sigma^*$ \\
$a \vee b$              & $\{a,b,c\}$ & $\{a,b,c\}$ & $\partial_c(a\vee b)=\varnothing$, still non-$\Sigma^*$ \\
$a^*$                   & $\{a,b\}$   & $\{a,b\}$  & $\partial_a(a^*)=a^*$, non-$\Sigma^*$; $b$ starts $[b]$ \\
$\neg a$ (double neg.)  & $\{a,b\}$   & $[a]$      & $\partial_b(\neg a)=\Sigma^*$ $\Rightarrow$ $b$ excluded$^\dagger$ \\
$a \cdot \Sigma^*$      & $\{a,b\}$   & $[b]$      & $\partial_a(a\cdot\Sigma^*)=\Sigma^*$ $\Rightarrow$ $a$ excluded$^\dagger$ \\
$a\cdot\Sigma^*\vee b\cdot\Sigma^*$ & $\{a,b,c\}$ & $[c]$ & both $a,b$ excluded; only $c$ avoids$^\dagger$ \\
$\neg a \wedge \neg b$  & $\{a,b,c\}$ & $\{a,b\}$  & De Morgan: $=\neg(a\vee b)$; $\partial_c$ is $\Sigma^*$$^\dagger$ \\
\bottomrule
\end{tabular}
\end{table}


\paragraph{Two worked cases.}
Two small examples fix the intuition for Equation~\eqref{eq:first-not}.
When $r = a\cdot\Sigma^*$ and $\Sigma_0 = \{a,b\}$, the naive rule
would return $\{a\}$, but $L(\neg r)$ consists of words beginning with
events other than $a$: $\partial_a(a\cdot\Sigma^*)$ normalises to $\univ$,
excluding $a$, while $\partial_b(a\cdot\Sigma^*) = \varnothing \neq \univ$
includes $b$, giving the correct $\{b\}$.
When $r = \varepsilon$ and $\Sigma_0 = \{a,b\}$, the naive rule returns
$\varnothing$ (since $\mathit{first}(\varepsilon) = \varnothing$), but
$\neg\varepsilon$ accepts all non-empty traces: $\partial_a(\varepsilon) =
\varnothing \neq \univ$ includes $a$, and likewise $b$, giving $\{a,b\}$.
Table~\ref{tab:firstwith-not} extends these to seven further patterns.
The first four rows confirm that complement does not exclude any alphabet
event when the derivative is non-total.
The double-negation row shows correct restoration of the original first set:
$\partial_b(\neg a) = \Sigma^*$ is total, so $b$ is filtered out.
In $\neg(a \cdot \Sigma^*)$, $\partial_a(a \cdot \Sigma^*)$ normalises to $\Sigma^*$,
excluding $a$; only $b$ survives.
Similarly, $\neg(a\cdot\Sigma^* \vee b\cdot\Sigma^*)$ excludes both $a$ and $b$,
leaving only $c$, while the De Morgan row verifies that
$\neg(\neg a \wedge \neg b) = a \vee b$ yields $\{a,b\}$.


\subsection{Embedding LTL\texorpdfstring{\textsubscript{f}}{f}}
\label{sec:ltl}

Because complement is a first-class constructor, the standard
LTL\textsubscript{f} operators translate directly into RE, with no
automaton construction needed.
Table~\ref{tab:ltl} gives the full mapping.

\begin{table}[ht]
\caption{LTL\textsubscript{f} to RE translation.
         $\Sigma^1 \triangleq \mathtt{Single\;Wildcard}$,
         $\univ \triangleq \neg\varnothing$.}
\label{tab:ltl}
\begin{tabular}{ll}
\toprule
LTL formula & RE $\sem{\cdot}$ \\
\midrule
$\top$              & $\univ$ \\
$\bot$              & $\varnothing$ \\
$e$                 & $\{e\}$ \\
$\neg\phi$          & $\neg\sem{\phi}$ \\
$\phi \wedge \psi$  & $\sem{\phi} \cap \sem{\psi}$ \\
$\phi \vee \psi$    & $\sem{\phi} \cup \sem{\psi}$ \\
$X\,\phi$           & $\Sigma^1 \cdot \sem{\phi}$ \\
$F\,\phi$           & $\univ \cdot \sem{\phi}$ \\
$G\,\phi$           & $\neg(\univ \cdot \neg\sem{\phi})$ \\
$\phi\,U\,\psi$     & $\mathit{step}(\phi)^* \cdot \sem{\psi}$
                      \quad ($\phi$ propositional) \\
\bottomrule
\end{tabular}
\end{table}

The complement operator does the heavy lifting:
$G\,\phi \triangleq \neg F\neg\phi$ unfolds to $\neg(\univ\cdot\neg\sem{\phi})$
using two applications of $\neg$; $F\,\phi$ and $X\,\phi$ need only concatenation.
The $U$ case uses $\mathit{step}(\phi)$, the length-1 RE for a single event
satisfying $\phi$, which is defined only when $\phi$ is propositional.

%% =========================================================================
\section{The \texttt{Pledge} Monad}
\label{sec:effectful}
%% =========================================================================

\subsection{The \texttt{Composable} Class}

We abstract the specification algebra via a type class:

\begin{lstlisting}
class Composable a where
  concatenation :: a -> a -> a   -- (<>), unit: empty
  conjunction   :: a -> a -> a   -- (/\), unit: universe
  empty         :: a             -- identity for <>  (varepsilon)
  universe      :: a             -- identity for /\  (neg-emptyset)
  subtraction   :: a -> a -> a   -- (\)
\end{lstlisting}


The five operations form exactly the interface required by the bind formula.
Here \lstinline{concatenation} sequences effects; \lstinline{conjunction}
combines simultaneous constraints.
The identities \lstinline{empty} and \lstinline{universe} appear in
\lstinline{pure}. In particular, the operation \lstinline{subtraction r1 r2} computes the residual of the
\emph{second} argument with respect to the first: it asks ``which part of
the obligation $r_2$ is not yet discharged, given that a trace matching
$r_1$ has been observed?''
The first argument $r_1$ is the \emph{observed} trace (a postcondition);
the second argument $r_2$ is the \emph{pending obligation} (a precondition
or future condition to be discharged).
Two base cases fix the endpoints: $\varepsilon \setminus r_2 = r_2$
(nothing observed, the full obligation $r_2$ remains) and
$\Sigma^* \setminus r_2 = \varepsilon$ (everything observed, the obligation
is fully discharged).
The general case advances both $r_1$ and $r_2$ in lockstep: for each
event $e$ that $r_1$ can begin with, the residual after $e$ is
$\partial_e(r_1) \setminus \partial_e(r_2)$, consuming one event from
each side simultaneously.

\emph{Quick example.}
Let the postcondition $r_1 = \mathtt{putMVar}(v)$ (\lstinline{putMVar(v)}
just occurred) and the obligation $r_2 = \finally{\mathtt{putMVar}(v)}$
(\lstinline{putMVar(v)} must eventually occur).
The event $\mathtt{putMVar}(v)$ in $r_1$ satisfies $\finally{\mathtt{putMVar}(v)}$,
so the residual $r_1 \setminus r_2 = \univ$ and the obligation is fully
discharged.
Conversely, if $r_1 = \mathtt{send}(v)$ instead, then
$\mathtt{send}(v) \setminus \finally{\mathtt{putMVar}(v)} =
\finally{\mathtt{putMVar}(v)}$: a \lstinline{send} event makes no progress
toward the \lstinline{putMVar} obligation, leaving it intact.

We write $\setminus$ for \lstinline{subtraction} uniformly throughout the
paper and in all listings; the library itself exposes this as a two-argument
top-level function rather than as an infix Haskell operator, so no clash
with reserved syntax arises.

\subsection{The \texttt{Pledge} Type}
\label{sec:future-cond}

The \lstinline{future} field is \emph{data-dependent}: it is a function of
the return value \lstinline{a}, allowing obligations to refer to the specific
resource handle that an operation produces.
The helper \lstinline{evalFuture} evaluates this function at the
computation's own return value.

\begin{lstlisting}
data Pledge eff a = Pledge
  { ret    :: a
  , pre    :: eff        -- what must have occurred before
  , post   :: eff        -- what this computation produces
  , future :: a -> eff   -- obligation indexed by the return value
  }

evalFuture :: Pledge eff a -> eff
evalFuture e = future e (ret e)
\end{lstlisting}

We call the value field \lstinline{ret} rather than the more common
\lstinline{val} specifically to mirror \lstinline{Monad}'s \lstinline{return}
and the mathematical presentation below; we use \lstinline{ret} only as a
record field name and never define a function of that name, so it does not
collide with monadic \lstinline{return}/\lstinline{pure}.

\subsection{Monad Instances}

The \lstinline{pure} function introduces a value with no effect and no obligation:

\begin{lstlisting}
pure x = Pledge { ret = x, pre = universe, post = empty, future = \_ -> universe }
\end{lstlisting}

The bind operator sequences two computations; write $e_2 = f(\ret(e))$.

\begin{lstlisting}
e >>= f = let fe = f (ret e) in Pledge
  { ret    = ret fe
  , pre    = pre e /\ (post e \\ pre fe)
  , post   = post e <> post fe
  , future = \_ -> (post fe \\ future e (ret e)) /\ future fe (ret fe)
  }
\end{lstlisting}

\paragraph{Future condition.}
Each \lstinline{future} function is applied to the corresponding
\lstinline{ret} before the obligation is propagated:

\begin{equation}
  \fut(e \mathbin{>\!\!>\!\!=} f)
  \;=\;
  \bigl(\post(e_2) \setminus \fut(e)(\ret(e))\bigr)
  \;\wedge\;
  \fut(e_2)(\ret(e_2))
  \label{eq:future-bind}
\end{equation}

This removes from $e$'s future obligation (evaluated at $e$'s return value)
those traces that $e_2$'s postcondition satisfies, then intersects the
residual with $e_2$'s own future obligation (evaluated at $e_2$'s return
value).
When the result normalises to $\univ$, all obligations are discharged.

\paragraph{Precondition.}

\begin{equation}
  \pre(e \mathbin{>\!\!>\!\!=} f)
  \;=\;
  \pre(e) \;\wedge\; \bigl(\post(e) \setminus \pre(e_2)\bigr)
  \label{eq:pre-bind}
\end{equation}

The term $\post(e) \setminus \pre(e_2)$ is the \emph{residual
precondition}: the part of $e_2$'s precondition not already discharged
by $e$'s postcondition.
Both must hold simultaneously in the same input history.

When $\post(e)$ does not cover $\pre(e_2)$, the residual is $\varnothing$,
flagging a violation.
When it fully covers $\pre(e_2)$, the residual is $\varepsilon$ and the
combined precondition reduces to $\pre(e)$ (the standard Hoare
sequencing rule).
We flag immediately, as a pointer forward to \S\ref{sec:monad-laws}, that
Equation~\eqref{eq:pre-bind} is exactly the place where \lstinline{Pledge}
departs from being a monad in the fully-generic, four-field sense: the
\lstinline{pre} field is carried alongside the monadic triple
$(\ret,\post,\fut)$ as a Hoare-style contract annotation, and the three
monad laws of \S\ref{sec:monad-laws} are stated for that triple only.

\subsection{The Applicative Instance}

The \lstinline{Pledge} monad also satisfies the \lstinline{Applicative}
class, which provides \lstinline{<*>} for combining a function-valued
computation with an argument-valued computation \emph{without}
data-dependency between them.
Writing $e_f$ for the function computation and $e_x$ for the argument
computation:

\begin{lstlisting}
ef <*> ex = Pledge
  { ret    = ret ef (ret ex)
  , pre    = pre ef /\ (post ef \\ pre ex)
  , post   = post ef <> post ex
  , future = \_ -> (post ex \\ future ef (ret ef)) /\ future ex (ret ex)
  }
\end{lstlisting}

The structure mirrors \lstinline{>>=}, i.e., preconditions and postconditions
compose in the same way, and the future condition is the conjunction of
the two evaluated futures after subtracting the continuation's
postcondition from the first computation's future.
The key difference is that \lstinline{<*>} cannot make $e_x$
\emph{depend} on the return value of $e_f$; \lstinline{>>=} passes
$\ret(e)$ into $f$, enabling the data-dependent future conditions
exemplified by \lstinline{tlsHandshake} and \lstinline{takeMVar}.
We define \lstinline{<*>} independently, by direct analogy with
\lstinline{>>=}, rather than deriving it via \lstinline{ap} from
\lstinline{Monad}; doing so lets the \lstinline{Applicative} instance be
given, and reasoned about, without committing to the full \lstinline{Monad}
machinery, which matters for the semiring-weighted instance of
\S\ref{sec:wre-instance} where some obligation domains admit sequencing
without needing data-dependent binding.

\subsection{Effects: How the Specification Layer Meets Real Code}
\label{sec:effects}

The presentation so far --- and every \lstinline{Composable} instance in
this paper --- treats \lstinline{eff} as a pure data type: \lstinline{RE},
\lstinline{GuardedRE}, \lstinline{WRE w}, or \lstinline{SL}. Correspondingly,
every annotated primitive we define (\lstinline{tlsHandshake},
\lstinline{takeMVar}, \lstinline{malloc}, \lstinline{alloc}, \ldots) is a
\emph{pure} \lstinline{Pledge} value: constructing it performs no \lstinline{IO},
touches no \lstinline{MVar}, and opens no file. This is deliberate, and it is
worth being explicit about why, and about what it means for the claim that
Pledge ``integrates with the \lstinline{mtl} stack.''

\paragraph{Why \lstinline{eff} is always pure.} A \lstinline{Pledge eff a}
value is a specification, not an implementation. Its three annotation
fields describe a trace, a heap shape, or a weight --- mathematical objects
that \lstinline{subtraction} can manipulate algebraically. If \lstinline{eff}
itself were \lstinline{IO} or contained live resources, \lstinline{subtraction}
would need to run and compare real IO actions, which is neither decidable
nor the point: the value of the derivative-based algebra of
\S\ref{sec:deriv} is precisely that it operates on syntax, without executing
anything. Keeping \lstinline{eff} pure is what makes residual annotations
computable at all.

\paragraph{What ``integrates with \texttt{mtl}'' means concretely.}
\lstinline{Pledge eff} is an ordinary \lstinline{Monad} instance, so it can be
used with plain \lstinline{do}-notation, \lstinline{sequence}, and the other
combinators any Haskell programmer already knows; no bespoke notation or
\lstinline{RebindableSyntax} is required, unlike the indexed-monad approaches
discussed in \S\ref{sec:related}. But we do \emph{not} claim that a
\lstinline{Pledge} value replaces or transparently wraps an \lstinline{IO},
\lstinline{StateT}, or \lstinline{ReaderT} computation, and no example in this
paper runs \lstinline{Pledge} and real effects in the same monad transformer
stack. We revise the earlier, overstated phrasing of the abstract
accordingly; what we mean is the narrower, and we believe still useful,
claim that a \lstinline{Pledge} specification can be constructed and composed
using the same \lstinline{Monad}/\lstinline{Applicative} vocabulary that
\lstinline{mtl}-style code already uses, so a Haskell programmer does not
need to learn a second composition discipline to read or write one.

\paragraph{The intended integration pattern: a shadow specification.}
Concretely, we envisage \lstinline{Pledge} used as a \emph{shadow}
specification that runs alongside, but separately from, the real effectful
code, checked once the specification term is fully built:

\begin{lstlisting}
main = do
  let spec = tlsHandshake ctx >>= \c -> tlsSend c >> tlsBye c   -- pure Pledge RE term
  assert (normalize (pre        spec) == top) "precondition violated"
  assert (normalize (evalFuture spec) == top) "future obligation unmet"
  realCtx <- Network.TLS.handshake realParams   -- the actual IO effect
  ...
\end{lstlisting}

\noindent
The \lstinline{spec} value and the real \lstinline{IO} computation are
constructed from the same call sequence by convention, and their
correspondence is the programmer's responsibility, not something Pledge
checks; keeping them in sync as the real code evolves is exactly the cost of
this shadow-specification approach, and we discuss tighter integration
(deriving the shadow automatically from an effect-handler row, so the two
cannot drift) as future work in \S\ref{sec:discussion}.

\paragraph{On the type of \texttt{future}.}
A closely related question is why \lstinline{future} has type
\lstinline{a -> eff} rather than simply \lstinline{eff}, given that every
example in \S\ref{sec:examples} applies \lstinline{future} only to
\lstinline{ret} (via \lstinline{evalFuture}), never to an arbitrary
argument. The answer is data-dependency at the point where \lstinline{future}
is \emph{written}, not where it is applied: \lstinline{takeMVar v} does not
know its own return value when its \lstinline{Pledge} record is constructed
in isolation, so its future condition must be expressed as a function
\lstinline{\char`\\r -> finally (Atom "putMVar" [Num r])} that names ``whichever
handle this operation returns.'' The function type only becomes visible as a
function inside \lstinline{>>=}, where the record is applied to the
\emph{corresponding} \lstinline{ret}, consistently, at both call sites in
Equation~\eqref{eq:future-bind}. We chose to make \lstinline{post} a plain
\lstinline{eff} rather than a function of the same shape because no example
in this paper requires a postcondition to depend on the value it is itself
reporting; a symmetric design in which \lstinline{post} is also
\lstinline{a -> eff}, applied to the previous computation's return value, is
a reasonable and strictly more general alternative that we did not need for
these case studies but do not rule out.

%% =========================================================================
\section{The RE Instance}
\label{sec:re-instance}
%% =========================================================================

The \lstinline{Composable RE} instance maps the algebra operations onto
\lstinline{RE} constructors and implements subtraction via derivatives:

\begin{lstlisting}
instance Composable RE where
  concatenation = Seq
  conjunction   = And
  empty         = Epsilon
  universe      = anything      -- = Not Bot

  subtraction Epsilon   r2 = r2
  subtraction (Not Bot) _  = Epsilon   -- Sigma* \\ r = eps (trivially satisfied)
  subtraction r1 r2 =
    let alph   = atoms r1 `union` atoms r2   -- combined alphabet
        evts   = firstWith alph r1
        step e = subtraction (normalize (derivative e r1))
                              (normalize (derivative e r2))
    in foldr Or Bot (map step evts)

normalize :: RE -> RE   -- RE-specific; not part of Composable
normalize r = {- ... algebraic simplification ... -}
\end{lstlisting}

The second clause handles the universal case:
$\univ \setminus r = \varepsilon$ for any $r$.
The general clause collects the alphabet from \emph{both} operands
before calling \lstinline{firstWith}, ensuring events appearing only in
$r_2$ are still considered under complement via Equation~\eqref{eq:first-not}.
The \lstinline{normalize} function is defined outside the instance and
applies the RE-specific laws listed in Table~\ref{tab:normalize-laws},
which exploit the complement constructor.

\paragraph{Termination.}
The recursive \lstinline{subtraction} call terminates because
\lstinline{normalize} reduces every RE to a canonical representative and,
for any regular expression $r$, the set of distinct normalised derivatives
$\{\,\lceil\partial_{w}(r)\rceil \mid w \in \Sigma^*\,\}$ is \emph{finite}
--- a classical result due to Brzozowski~\cite{brzozowski1964}, refined to
include complement by Owens et al.~\cite{owens2009}.
Each recursive call to \lstinline{subtraction} is made with the
normalised derivative of the \emph{previous} operands, so the recursion
explores a path in this finite set and must eventually reach a base case
($\varepsilon$ or $\univ$) or a previously visited pair.

\begin{table}[ht]
\caption{Normalisation laws applied by \lstinline{normalize}.
         All rules are applied recursively bottom-up.}
\label{tab:normalize-laws}
\begin{tabular}{ll}
\toprule
Law & Constructor \\
\midrule
$\varepsilon \cdot r = r$, $r \cdot \varepsilon = r$ & \lstinline{Seq} \\
$\varnothing \vee r = r$, $r \vee \univ = \univ$ & \lstinline{Or} \\
$\varnothing \wedge r = \varnothing$, $r \wedge \univ = r$ & \lstinline{And} \\
$\neg\neg r = r$ & \lstinline{Not} \\
$\neg(r_1\vee r_2) = \neg r_1\wedge\neg r_2$ & \lstinline{Not} \\
$\neg(r_1\wedge r_2) = \neg r_1\vee\neg r_2$ & \lstinline{Not} \\
$\neg\varnothing = \univ$, $\neg\univ = \varnothing$ & \lstinline{Not} \\
$\varnothing^* = \varepsilon$, $\varepsilon^* = \varepsilon$ & \lstinline{Star} \\
\bottomrule
\end{tabular}
\end{table}

\subsection{Soundness}
\label{sec:soundness}

The examples of \S\ref{sec:examples} read a fully-discharged residual
($\univ$ in \lstinline{future}, $\univ$ in \lstinline{pre}) as a certificate
that the composed specification is contract-correct. We now state and prove
the theorem that licenses this reading for the $\RE$ instance, addressing
the soundness gap raised by both Reviewer~A and the program committee.

For an $\RE$ value $r$ and a finite trace $w = a_1 a_2 \cdots a_n$ over
events, write $w \models r$ for $w \in L(r)$, the standard language
membership relation, and $\partial_w(r)$ for the iterated derivative
$\partial_{a_n}(\cdots\partial_{a_1}(r)\cdots)$.

\begin{lemma}[Derivative correctness]
\label{lem:deriv-correct}
For all $r$ and all traces $w$, $w \models r$ if and only if $\nu(\partial_w(r)) = \mathtt{True}$.
\end{lemma}
\begin{proof}
By induction on $w$. The base case $w = \varepsilon$ is exactly the
definition of $\nu$. For $w = a \cdot w'$, $w \models r$ iff
$w' \models \partial_a(r)$, which is the standard derivative correctness
property~\cite{brzozowski1964}, extended to \lstinline{Not} by
Equation~\eqref{eq:deriv-not} and to \lstinline{And} by the direct clause
in \S\ref{sec:re}; the inductive hypothesis on $w'$ and $\partial_a(r)$
then gives $\nu(\partial_{w'}(\partial_a(r))) = \nu(\partial_w(r))$, as
required.
\end{proof}

\begin{lemma}[Subtraction correctness]
\label{lem:sub-correct}
For all $r_1, r_2$ and all traces $w$: if $w \models r_1$, then
$w \models r_2$ if and only if $\varepsilon \models \partial_w(r_1) \setminus \partial_w(r_2)$
is witnessed along the recursion, i.e.\ $\nu(\partial_w(r_1) \setminus \partial_w(r_2))$
reduces, via the base cases of \lstinline{subtraction}, to \lstinline{True}
exactly when $\nu(\partial_w(r_2)) = \mathtt{True}$.
\end{lemma}
\begin{proof}
By induction on the structure of the \lstinline{subtraction} recursion. The
base case $\mathtt{subtraction}\;\varepsilon\;r_2 = r_2$ applies when $r_1$'s
derivative has been driven to $\varepsilon$, i.e.\ $w$ is exactly a word of
$L(r_1)$ with nothing left to consume; nullability of the residual is then
nullability of $r_2$ itself, matching Lemma~\ref{lem:deriv-correct}. The base
case $\mathtt{subtraction}\;(\neg\varnothing)\;r_2 = \varepsilon$ applies when
$r_1$'s derivative has become the universal language, in which case every
continuation is already in $L(r_1)$, so the residual should record ``fully
observed,'' i.e.\ nullable unconditionally, which $\varepsilon$ is. The
recursive case advances $r_1$ and $r_2$ by the same event $e$ and applies
the inductive hypothesis at $\partial_e(r_1)$, $\partial_e(r_2)$.
\end{proof}

\begin{theorem}[Soundness of $\RE$-residual checking]
\label{thm:soundness}
Let $e :: \mathtt{Pledge}\;\RE\;a$ be a specification term built from
\lstinline{pure} and \lstinline{>>=} applications of $\RE$-annotated
primitives, and let $w$ be any finite trace of events actually emitted by a
terminating execution whose call structure matches $e$ (i.e.\ $w$ is a
concatenation of the \lstinline{post} events of the primitives invoked, in
invocation order). If $\mathtt{normalize}(\mathtt{evalFuture}\;e) = \univ$
and $\mathtt{normalize}(\mathtt{pre}\;e) = \univ$, then $w$ satisfies every
future condition and every precondition named by the primitives composing
$e$.
\end{theorem}
\begin{proof}[Proof sketch]
By induction on the structure of $e$. For $e = \mathtt{pure}\;x$, both
residuals are $\univ$ by definition and there are no obligations to satisfy,
so the claim holds vacuously. For $e = e_1 \mathbin{>\!\!>\!\!=} f$, the
residual future of $e$ is given by Equation~\eqref{eq:future-bind}, which by
Lemma~\ref{lem:sub-correct} is $\univ$ exactly when the trace segment
corresponding to $f(\ret(e_1))$'s postcondition discharges $e_1$'s future
obligation (evaluated at $\ret(e_1)$) and $f(\ret(e_1))$'s own future
obligation is separately $\univ$; the inductive hypothesis applied to $e_1$
and to $f(\ret(e_1))$ then gives that the corresponding trace segments
satisfy their respective obligations, and concatenating these segments (by
the definition of $w$) gives that $w$ satisfies the future obligation of $e$.
The precondition case is symmetric, using Equation~\eqref{eq:pre-bind} and
Lemma~\ref{lem:sub-correct} in place of Equation~\eqref{eq:future-bind}.
\end{proof}

\noindent
Theorem~\ref{thm:soundness} is stated for the $\RE$ instance because its
proof leans on the derivative correctness of Lemma~\ref{lem:deriv-correct},
which is specific to regular languages. The $\GRE$ instance inherits this
soundness for its RE component and additionally requires the Presburger
side to be checked by a sound decision procedure (Z3), which we take as
given rather than re-proving. The $\WRE$ instance admits an analogous
statement with $\nu$ replaced by \lstinline{wNullable} and the conclusion
weakened to a numeric bound on the accumulated weight, which we state but
do not prove in detail. The $\SL$ instance, as discussed in
\S\ref{sec:sl-instance}, does not admit an analogous soundness theorem as it
stands, precisely because its \lstinline{subtraction} records the wand
symbolically without a semantic model of the heap; establishing one is
exactly the ``real solver'' extension identified as future work in
\S\ref{sec:discussion}.

%% =========================================================================
\section{The GuardedRE Instance}
\label{sec:guardedre-instance}
%% =========================================================================

The RE instance reasons about event traces but is entirely silent about
arithmetic state.
The $\GRE$ instance closes this gap: a guarded regular expression pairs
a trace RE with a Presburger arithmetic predicate over a concrete heap,
so that a single annotation can simultaneously enforce a protocol ordering
\emph{and} an arithmetic bound.

\subsection{The \texttt{GuardedRE} Type}

\begin{lstlisting}
data GuardedRE = GuardedRE PPred RE
    deriving (Eq)
\end{lstlisting}

A state $(\text{heap},\text{trace})$ satisfies \lstinline{GuardedRE p r} iff
\[
  \text{heap} \models p \quad\wedge\quad \text{trace} \in L(r).
\]
The two sides are independent: $p$ is a static constraint on the heap,
while $r$ advances through Brzozowski derivatives as events arrive.

Two smart constructors lift a pure RE or a pure Presburger predicate:

\begin{lstlisting}
fromRE    :: RE    -> GuardedRE
fromRE r   = GuardedRE PTrue r        -- no heap constraint

fromPPred :: PPred -> GuardedRE
fromPPred p = GuardedRE p universe    -- accept any trace
\end{lstlisting}

\subsection{The \texttt{Composable GuardedRE} Instance}

\begin{lstlisting}
instance Composable GuardedRE where
    concatenation (GuardedRE p1 r1) (GuardedRE p2 r2) =
        GuardedRE (normalizePPred (PAnd p1 p2)) (normalize (Seq r1 r2))
    conjunction (GuardedRE p1 r1) (GuardedRE p2 r2) =
        GuardedRE (normalizePPred (PAnd p1 p2)) (normalize (And r1 r2))
    empty    = GuardedRE PTrue Epsilon
    universe = GuardedRE PTrue (Not Bot)   -- = anything
    subtraction (GuardedRE p1 r1) (GuardedRE p2 r2) =
        GuardedRE (normalizePPred (PAnd p1 p2))
                  (normalize (reSubtraction r1 r2))
\end{lstlisting}

\noindent
Both \lstinline{concatenation} and \lstinline{subtraction} conjoin
both heap predicates: the residual must respect the constraints from
both operands.
The RE component uses the same \lstinline{reSubtraction} and
\lstinline{normalize} as the plain RE instance (named
\lstinline{reSubtraction} here purely to distinguish it, in this listing,
from the outer \lstinline{GuardedRE} \lstinline{subtraction} being
defined; it is the identical function called \lstinline{subtraction} in
\S\ref{sec:re-instance}).

\subsection{Presburger Expressions and Normalisation}

Presburger predicates are built over a small arithmetic expression
language:

\begin{lstlisting}
data PExpr = Lit Int | ValAt Addr | PAdd PExpr PExpr | PSub PExpr PExpr

data PPred
    = PTrue | PFalse      -- canonical true/false
    | PLt PExpr PExpr | PLe PExpr PExpr | PEq PExpr PExpr | PGe PExpr PExpr
    | PNot PPred | PAnd PPred PPred
\end{lstlisting}

\noindent
\lstinline{Lit k} is the integer literal $k$; \lstinline{ValAt a} reads the
current value stored at heap address \lstinline{a} (so, e.g.,
\lstinline{PLt (ValAt a) (Lit maxVal)} is the predicate $h[a] < \mathit{maxVal}$
used in \S\ref{sec:ex-bounded}); \lstinline{PAdd}/\lstinline{PSub} build
compound arithmetic expressions from these.

Composing multiple \lstinline{GuardedRE} values accumulates a tree of
\lstinline{PAnd} nodes.
The \lstinline{normalizePPred} function flattens this tree, removes
duplicates, substitutes equalities (e.g.\ $h[a]=k \Rightarrow$ replace
$\mathtt{ValAt}\;a$ by $\mathtt{Lit}\;k$ throughout), evaluates literal
comparisons, and replaces $\lnot\mathbf{true}$ with the canonical
\lstinline{false} constructor \lstinline{PFalse}:

\begin{lstlisting}
normalizePPred :: PPred -> PPred
normalizePPred p = rebuildAnd (substEqs eqs atoms)
  where
    atoms = flattenAnd (normStep p)
    eqs   = [ (a, k) | PEq (ValAt a) (Lit k) <- atoms ]
\end{lstlisting}

The introduction of \lstinline{PFalse} prevents the solver from being
invoked with expressions such as $(\lnot\mathbf{true}) \wedge \cdots$ that
are syntactically unsatisfiable; these normalise immediately to
\lstinline{PFalse} without any solver call. Predicates that survive this
purely syntactic normalisation step are discharged by an SMT call to Z3
(via SBV) to determine satisfiability.

\subsection{Example: Bounded Counter}

A bounded counter stored at heap address $a$ with maximum value $M$
must follow the protocol $\mathtt{init}\cdot(\mathtt{inc}\mid\mathtt{dec})^*\cdot\mathtt{snapshot}$.
Increment is additionally guarded by $h[a] < M$ (no overflow); decrement
by $h[a] > 0$ (no underflow).

\begin{lstlisting}
initCounter :: Addr -> Pledge GuardedRE ()
initCounter a = Pledge
  { ret = ()
  , pre  = fromRE universe
  , post = GuardedRE (PEq (ValAt a) (Lit 0))
                     (Single (Atom "init" [Num a]))
  , future = \_ -> fromRE universe }

increment :: Addr -> Int -> Pledge GuardedRE ()
increment a maxVal = Pledge
  { ret = ()
  , pre  = GuardedRE (PLt (ValAt a) (Lit maxVal)) universe
  , post = fromRE (Single (Atom "inc" [Num a]))
  , future = \_ -> fromPPred (PGe (ValAt a) (Lit 0)) }
\end{lstlisting}

\noindent
The \lstinline{pre} of \lstinline{increment} carries
$h[a] < M$ (pure Presburger) conjoined with the universal RE.
After composing eleven increments against a bound of~10,
\lstinline{normalizePPred} derives the contradiction
$h[a] < 10 \wedge \lnot(h[a] < 10) \equiv \mathbf{false}$,
flagging the overflow before any execution takes place.

\paragraph{Bad programs and residuals.}
\begin{center}
\begin{tabular}{lll}
\toprule
Bad program & Residual \lstinline{pre} (Presburger) & Residual \lstinline{pre} (RE) \\
\midrule
\lstinline{overflow} (11 incs, max=10) & $\mathbf{false}$ & $\univ$ \\
\lstinline{underflow} (dec at zero)    & $\mathbf{false}$ & $\univ$ \\
\lstinline{noInit} (skip init)         & $\univ$           & $\varnothing$ \\
\bottomrule
\end{tabular}
\end{center}

%% =========================================================================
\section{The WeightedRE Instance}
\label{sec:wre-instance}
%% =========================================================================

The RE and $\GRE$ instances are Boolean: a future condition is either
satisfied or violated.
The $\WRE$ instance generalises membership to a \emph{weight} drawn from
a semiring, enabling probabilistic and quantitative obligations to be
expressed within the same \lstinline{Composable} framework.

\subsection{The Semiring Class}

\begin{lstlisting}
class (Eq w, Show w) => Semiring w where
    szero :: w          -- additive identity  (blocked path)
    sone  :: w          -- multiplicative identity
    sadd  :: w -> w -> w
    smul  :: w -> w -> w
\end{lstlisting}

\noindent
Setting $w = \mathtt{Bool}$ (with $\mathtt{sadd} = (\vee)$,
$\mathtt{smul} = (\wedge)$) recovers the Boolean RE instance.
Two further instances capture quantitative scenarios:

\begin{center}
\begin{tabular}{llll}
\toprule
Instance & \lstinline{szero} & \lstinline{sadd} & \lstinline{smul} \\
\midrule
\lstinline{Prob}     (probability) & $0$        & $\min(p_1+p_2,\,1)$ & $p_1 \times p_2$ \\
\lstinline{Tropical} (min-cost)    & $+\infty$  & $\min$               & $+$              \\
\bottomrule
\end{tabular}
\end{center}

\noindent
In the probability semiring, \lstinline{sadd} is bounded addition (sum of
probabilities capped at 1) and \lstinline{smul} is multiplication.
In the tropical semiring, \lstinline{sadd} is minimum (taking the
cheaper alternative) and \lstinline{smul} is addition (summing costs
along a path).

\subsection{The \texttt{WRE} Type and Operators}

\begin{lstlisting}
data WRE w
    = WBot               -- weight szero: no accepting path
    | WEps w             -- weight w: accept the empty trace
    | WSingle w Event    -- weight w: accept exactly this event
    | WSeq (WRE w) (WRE w)
    | WAdd (WRE w) (WRE w)
    | WAnd (WRE w) (WRE w)
    | WStar (WRE w)
\end{lstlisting}

\noindent
The generalised nullable function accumulates the weight of all
accepting paths through the empty trace:

\begin{lstlisting}
wNullable :: Semiring w => WRE w -> w
wNullable (WEps w)      = w
wNullable (WSeq r1 r2)  = smul (wNullable r1) (wNullable r2)
wNullable (WAdd r1 r2)  = sadd (wNullable r1) (wNullable r2)
wNullable (WAnd r1 r2)  = smul (wNullable r1) (wNullable r2)
wNullable (WStar r)     = sone
wNullable _             = szero
\end{lstlisting}

\noindent
The functions \lstinline{wDerivative} and \lstinline{wSubtraction} follow
the same structure as their Boolean counterparts, with semiring operations
substituted for Boolean ones.
The \lstinline{Composable (WRE w)} instance (enabled by
\lstinline{FlexibleInstances}) maps \lstinline{concatenation} to
\lstinline{WSeq}, \lstinline{conjunction} to \lstinline{WAnd}, and
\lstinline{subtraction} to \lstinline{wSubtraction}.

\subsection{Derived Forms}

\begin{lstlisting}
wTop     :: Semiring w => WRE w
wTop      = WStar (WSingle sone Wildcard)   -- sone-weighted anything

wFinally :: Semiring w => w -> Event -> WRE w
wFinally weight e = WSeq wTop (WSingle weight e)

wGlobally :: Semiring w => w -> Event -> WRE w
wGlobally weight e = WStar (WSingle weight e)
\end{lstlisting}

\subsection{Example: Probabilistic Memory (Prob Semiring)}

In a probabilistic setting, each \lstinline{free(a)} has an associated
probability of being called.
Each \lstinline{malloc} call carries a future condition asserting that
\lstinline{free(a)} must be called with probability at least $p$:

\begin{lstlisting}
malloc :: Addr -> Double -> Pledge (WRE Prob) Addr
malloc a p = Pledge
  { ret = a, pre = wTop
  , post   = WSingle sone (Atom "malloc" [Num a])
  , future = \r -> wFinally (Prob p) (Atom "free" [Num r]) }
\end{lstlisting}

The residual weight of \lstinline{wNullable} at the end of a program
measures how much of the obligation has been satisfied:
a weight of \lstinline{Prob 1.0} indicates that the obligation is fully
discharged; a weight strictly less than~$1$ quantifies the probability shortfall.
A call \lstinline{malloc a 0.9} whose corresponding \lstinline{free a} is
never composed into the program leaves a residual weight of exactly
\lstinline{Prob 0.9}: the obligation is not a contract the caller must
write, since the caller cannot control the run's probabilistic outcome; it
is the declared quantitative claim itself, and the residual reports how much
of that claim the composed call sequence still owes, which is the reading
we intend when we describe $\WRE$ obligations as ``an output of an
analysis'' in \S\ref{sec:intro}.

\subsection{Example: Min-Cost Scheduling (Tropical Semiring)}

Under the tropical semiring, \lstinline{WSingle (Tropical c) e}
assigns cost $c$ to emitting event $e$; \lstinline{wNullable} then
returns the minimum total cost of reaching an accepting state:

\begin{lstlisting}
submitJob :: Int -> Double -> Pledge (WRE Tropical) ()
submitJob jobId cost = Pledge
  { ret = ()
  , pre  = wTop
  , post = WSingle (Tropical cost) (Atom "submit" [Num jobId])
  , future = \_ ->
      wFinally (Tropical 0) (Atom "complete" [Num jobId]) }
\end{lstlisting}

Composing several \lstinline{submitJob} calls yields a future whose
\lstinline{wNullable} is the minimum total cost of completing all
submitted jobs.
A program that omits a \lstinline{complete} event for some job~$i$
retains \lstinline{wFinally (Tropical 0) (complete i)} in the
residual, and the overall weight is $+\infty$ (i.e.\ \lstinline{szero}),
flagging the missing completion.

\paragraph{Absence of weighted complement.}
The \lstinline{WRE} type does not include a \lstinline{WNot} constructor.
This is a deliberate restriction: complement over a semiring weight does
not have a canonical definition, because Boolean negation ($0 \leftrightarrow 1$)
does not generalise uniformly to an arbitrary semiring
(e.g.\ what is the ``negation'' of a probability $0.7$, or of a tropical
cost $5$?).
As a consequence, the LTL$_f$ embedding from Section~\ref{sec:ltl}
does not transfer to the \lstinline{WRE} instance:
operators that rely on complement --- \lstinline{G}, \lstinline{U},
and \lstinline{never} --- are unavailable.
The operators \lstinline{wFinally} and \lstinline{wGlobally} above are
defined directly via \lstinline{WStar} and \lstinline{WSeq} without
complement, and cover the common quantitative patterns. This is the
concrete sense, referenced in \S\ref{sec:intro}, in which $\WRE$ is fully
usable for the obligations it targets while giving up part of the $\RE$
instance's generality.

%% =========================================================================
\section{The SL Instance}
\label{sec:sl-instance}
%% =========================================================================

The \lstinline{Composable} typeclass is not inherently tied to
trace-based specifications.
We demonstrate its generality with a fourth instance over
\emph{separation-logic heap predicates}, where future conditions describe
what heap state must hold upon completion of the computation. As stated in
\S\ref{sec:intro}, this instance is included primarily to show that
\lstinline{Composable} generalises beyond traces; we are explicit
throughout this section, and again in \S\ref{sec:soundness}, about the
respects in which it falls short of the RE instance's guarantees.

\subsection{Structural Parallel with RE}

Table~\ref{tab:sl-re-parallel} shows the role each SL construct plays in
the \lstinline{Composable} interface.

\begin{table}[ht]
\caption{Structural parallel between the four \lstinline{Composable} instances.}
\label{tab:sl-re-parallel}
\begin{tabular}{lllll}
\toprule
Operation & $\RE$ & $\GRE$ & $\WRE_w$ & $\SL$ \\
\midrule
\lstinline{concatenation} & $r_1 \cdot r_2$ & $(p_1\wedge p_2,\,r_1\cdot r_2)$ & \lstinline{WSeq} & $P \sep Q$ \\
\lstinline{conjunction}   & $r_1 \wedge r_2$ & $(p_1\wedge p_2,\,r_1\wedge r_2)$ & \lstinline{WAnd} & $P \wedge Q$ \\
\lstinline{empty}         & $\varepsilon$ & $(\mathbf{true},\varepsilon)$ & \lstinline{WEps sone} & $\mathbf{emp}$ \\
\lstinline{universe}      & $\univ$ & $(\mathbf{true},\univ)$ & \lstinline{wTop} & $\top$ \\
\lstinline{subtraction}   & Brzozowski quotient & quotient (RE) + $\wedge$ (PPred) & \lstinline{wSubtraction} & $P \wand Q$ \\
\bottomrule
\end{tabular}
\end{table}

\noindent
The bind formula is unchanged:
$\pre(e \mathbin{>\!\!>\!\!=} f) = \pre(e) \wedge (\post(e) \wand \pre(e_2))$
and
$\fut(e \mathbin{>\!\!>\!\!=} f) = (\post(e_2) \wand \fut(e)) \wedge \fut(e_2)$.
Here $\wand$ is the magic wand (separating implication):
$P \wand Q$ holds of heap $h$ iff for every heap $h'$ disjoint from $h$
satisfying $P$, the combined heap $h \cup h'$ satisfies $Q$.

\subsection{The SL Predicate Type}

\begin{lstlisting}
data SL
    = Emp              -- empty heap         (identity for SepStar)
    | Top              -- any heap           (identity for Conj)
    | Pure PPred       -- pure Presburger constraint (heap-independent)
    | Cell Addr Val    -- singleton: address a holds value v
    | SepStar SL SL    -- P * Q  separating conjunction
    | Conj SL SL       -- P /\ Q ordinary conjunction
    | Wand SL SL       -- P -* Q magic wand (residual)
    deriving (Eq, Show)
\end{lstlisting}

\subsection{Separating vs.\ Ordinary Conjunction}

$\mathtt{SepStar}\;P\;Q$ holds of heap $h$ iff $h$ splits into
disjoint parts $h_1 \uplus h_2$ with $P(h_1)$ and $Q(h_2)$; it is the
standard separating conjunction of ownership.
$\mathtt{Conj}\;P\;Q$ holds of $h$ iff the \emph{same} heap satisfies
both $P$ and $Q$.
Without pure constraints, \lstinline{Conj} on purely spatial predicates
is either trivial ($\top \wedge P = P$) or unsatisfiable
($\mathtt{Cell}\;l_1\;v_1 \wedge \mathtt{Cell}\;l_2\;v_2 = \mathbf{false}$
for $l_1 \neq l_2$).
The constructor becomes useful when one side is a pure constraint
$\lceil\phi\rceil \wedge (\ell \mapsto v)$,
combining a heap-independent condition with spatial ownership.

\subsection{Presburger Arithmetic for Pure Constraints}

Pure constraints reuse the same \lstinline{PPred}/\lstinline{PExpr}
language introduced in Section~\ref{sec:guardedre-instance}
(\lstinline{normalizePPred}, and the reading of \lstinline{ValAt a} as the
value currently stored at address \lstinline{a}).
The predicate \lstinline{Pure p} holds of any heap for which
\lstinline{evalPPred h p} is true; it enforces no ownership of cells.
The smart constructor \lstinline{pureSL} realises the standard SL box
notation: \lstinline{pureSL PTrue = Emp} and
\lstinline{pureSL p = Conj (Pure p) Emp} for non-trivial $p$.

\subsection{The \texorpdfstring{\lstinline{Composable SL}}{Composable SL} Instance}

\begin{lstlisting}
instance Composable SL where
    concatenation = SepStar
    conjunction   = Conj
    empty         = Emp
    universe      = Top

    subtraction Emp q = q           -- emp -* Q = Q
    subtraction Top _ = Emp         -- Top trivially covers any pre
    subtraction p   q = Wand p q    -- general: symbolic wand
\end{lstlisting}

\noindent
The two base cases mirror those of the RE instance:
$\varepsilon \setminus r = r$ becomes $\mathbf{emp} \wand Q = Q$
(providing no heap leaves the obligation unchanged), and
$\univ \setminus r = \varepsilon$ becomes $\top \wand Q = \mathbf{emp}$
(a $\top$ postcondition covers any precondition, leaving no residual).
The general case records the wand symbolically for later evaluation
or normalisation.

\subsection{Normalisation}

The \lstinline{normalizeSL} function applies the standard SL algebraic identities:

\begin{center}
\begin{tabular}{ll}
\toprule
Law & Constructor \\
\midrule
$\mathbf{emp} \sep Q = Q$, $P \sep \mathbf{emp} = P$ & \lstinline{SepStar} \\
$\top \wedge Q = Q$, $P \wedge \top = P$              & \lstinline{Conj} \\
$P \wedge P = P$ (idempotent)                          & \lstinline{Conj} \\
$\mathbf{emp} \wand Q = Q$                            & \lstinline{Wand} \\
$P \wand \top = \top$                                  & \lstinline{Wand} \\
\bottomrule
\end{tabular}
\end{center}


\paragraph{Symbolic recording, not semantic discharge.}
It is important to distinguish what the \lstinline{SL} instance does from
what it does not do.
The five normalisation rules above reduce \emph{syntactically} trivial
wands (e.g.\ $\mathbf{emp} \wand Q = Q$), but the general
\lstinline{Wand p q} case is recorded \emph{symbolically}: no heap model
is consulted, and no solver is invoked.
A residual such as $\mathtt{Wand}\;(\mathtt{Cell}\;0\;42)\;(\mathtt{Cell}\;0\;99)$
is retained as a data structure even though it is semantically
unsatisfiable ($0 \mapsto 42$ and $0 \mapsto 99$ cannot simultaneously
hold).
The \lstinline{SL} instance therefore \emph{records} obligations and
propagates them correctly through bind; it does not \emph{check} whether
the accumulated residual is satisfiable or entailed, and consequently
Theorem~\ref{thm:soundness}'s style of guarantee does not extend to it
(\S\ref{sec:soundness}).
Connecting residual \lstinline{SL} predicates to an external solver
(e.g.\ Z3 via the \lstinline{sbv} library, or a dedicated SL decision
procedure such as Smallfoot) is identified in Section~\ref{sec:discussion}
as the natural next engineering step.
All claims of ``violation detection'' for the SL instance in the
examples below should be read as ``the residual annotation is
non-trivially non-\lstinline{universe}'', which a programmer or tool can
then inspect further, exactly the caveat flagged for the SL row of
Table~\ref{tab:sl-re-parallel}.

%% =========================================================================
\section{Examples}
\label{sec:examples}
%% =========================================================================

We demonstrate \lstinline{Pledge} across use cases spanning all four instances.
Each example defines primitive operations annotated with
\lstinline{pre}, \lstinline{post}, and \lstinline{future},
composes them in the \lstinline{Pledge} monad, and inspects the
normalised annotations of the result.
A \lstinline{future} of $\univ$ (or \lstinline{wNullable = sone} for $\WRE$)
and a \lstinline{pre} of $\univ$ together certify temporal correctness
(Theorem~\ref{thm:soundness}, for $\RE$).
Where an obligation domain admits one, we favour examples whose
\lstinline{future} condition is genuinely non-trivial, since a
\lstinline{universe} future condition throughout a derivation would make no
use of the propagation machinery this paper is about.

\subsection{TLS Handshake Lifecycle}
\label{sec:ex-tls}

The \lstinline{tls} library obligation is representative of a broad class
of Haskell protocol obligations: a successful \lstinline{handshake} must
always be paired with a subsequent \lstinline{bye}.
We model the \lstinline{Context} as an integer identifier.

\begin{lstlisting}
type Context = Int

tlsHandshake :: Context -> Pledge RE Context
tlsHandshake ctx = Pledge
  { ret    = ctx
  , pre    = universe
  , post   = Single (Atom "handshake" [Num ctx])
  , future = \c -> finally (Atom "bye" [Num c]) }

tlsBye :: Context -> Pledge RE ()
tlsBye ctx = Pledge
  { ret    = ()
  , pre    = Single (Atom "handshake" [Num ctx])
  , post   = Single (Atom "bye" [Num ctx])
  , future = \_ -> universe }

tlsSend :: Context -> Pledge RE ()
tlsSend ctx = Pledge
  { ret    = ()
  , pre    = Single (Atom "handshake" [Num ctx])
  , post   = Single (Atom "send" [Num ctx])
  , future = \_ -> universe }
\end{lstlisting}

\begin{lstlisting}
-- Good: handshake followed by data exchange then bye
goodSession = do
  ctx <- tlsHandshake 1
  tlsSend ctx
  tlsBye ctx

-- Bad future: bye omitted; residual = finally(bye(1))
missingBye = do
  ctx <- tlsHandshake 1
  tlsSend ctx

-- Bad future: two contexts opened, only one closed
partialClose = do
  ctx1 <- tlsHandshake 1
  ctx2 <- tlsHandshake 2
  tlsBye ctx1
\end{lstlisting}

\noindent
The program \lstinline{missingBye} retains $\finally{\mathtt{bye}(1)}$ as
its residual future, naming the unclosed context precisely.
Meanwhile, \lstinline{partialClose} retains $\finally{\mathtt{bye}(2)}$:
the future condition is data-dependent, so the residual identifies
context~2 specifically, not a generic ``some context was not closed.''

\subsection{MVar Lifecycle}

The \lstinline{takeMVar} operation returns the taken handle; its
data-dependent future condition uses the lambda \lstinline{r} so
the obligation tracks the \emph{returned} handle.
The \lstinline{putMVar} operation carries as its future obligation
$\noUntil{\mathtt{putMVar}(v)}{\mathtt{takeMVar}(v)}$,
forbidding a second \lstinline{putMVar(v)} until
\lstinline{takeMVar(v)} resets the guard.
(The full definitions of \lstinline{takeMVar} and \lstinline{putMVar}
appear in Section~\ref{sec:future-cond}.) As discussed in
\S\ref{sec:effects}, \lstinline{MVar} here is the pure type synonym
\lstinline{type MVar = Int} used by the specification, standing in for a
real \lstinline{Control.Concurrent.MVar} handle; see \S\ref{sec:discussion}
for how a real program's \lstinline{MVar} usage would be linked to this
specification.

\begin{lstlisting}
-- Good: data-dependent  -  obligation discharged on the returned handle
takeAndPut n = do { v <- takeMVar n; putMVar v }

-- Good: re-acquire same MVar after releasing  -  noUntil is reset by takeMVar
takePutTakePut = do { v <- takeMVar 1; putMVar v; v <- takeMVar 1; putMVar v }

-- Bad future: putMVar(2) obligation remains (MVar 2 never returned)
missingPut = do { v1 <- takeMVar 1; putMVar v1; _ <- takeMVar 2; return () }

-- Bad pre: takeMVar(1) precondition not met
putWithoutTake = putMVar 1

-- Bad future: double-put immediately collapses future to Bot
doublePutImmediate = do { v <- takeMVar 1; putMVar v; putMVar v }

-- Bad future: double-put with intervening work, still collapses to Bot
doublePutWithWork = do
  { v1 <- takeMVar 1; v2 <- takeMVar 2; putMVar v1; putMVar v2; putMVar v1 }
\end{lstlisting}

The three examples below return non-unit values, showing that future
conditions are tracked regardless of the return type.

\begin{lstlisting}
-- Good: ret = MVar; future = finally(putMVar(1))  -  obligation visible in residual
takeReturnHandle :: Pledge RE MVar
takeReturnHandle = takeMVar 1

-- Bad future: ret = (MVar,MVar); future = finally(putMVar(1)) /\ finally(putMVar(2))
takeTwoReturnPair :: Pledge RE (MVar, MVar)
takeTwoReturnPair = do { v1 <- takeMVar 1; v2 <- takeMVar 2; return (v1, v2) }

-- Good: ret = MVar; future = noUntil(putMVar(1),takeMVar(1))  -  no pending finally
takePutReturnHandle :: Pledge RE MVar
takePutReturnHandle = do { v <- takeMVar 1; putMVar v; return v }
\end{lstlisting}

\subsection{Mutex / Lock Lifecycle}

\begin{lstlisting}
release :: Int -> Pledge RE ()
release m = Pledge
  { ret = ()
  , pre    = Single (Atom "acquire" [Num m])
  , post   = Single (Atom "release" [Num m])
  , future = universe }

-- Bad future: lock 1 not released
lockLeak = do { acquire 1; acquire 2; release 2 }

-- Bad pre: release without acquire
releaseWithoutAcquire = release 1
\end{lstlisting}

\subsection{Database Transactions}

The \lstinline{beginTx} operation carries a \emph{disjunctive} future
condition requiring that the transaction is eventually either committed
or rolled back; \lstinline{commit} and \lstinline{rollback} each
discharge this obligation:

\begin{lstlisting}
beginTx :: Pledge RE ()
beginTx = Pledge
  { ret = (), pre = universe
  , post   = Single (Atom "beginTx" [])
  , future = Or (finally (Atom "commit" []))
                (finally (Atom "rollback" [])) }

commit :: Pledge RE ()
commit = Pledge
  { ret = (), pre = Single (Atom "beginTx" [])
  , post   = Single (Atom "commit" [])
  , future = universe }

rollback :: Pledge RE ()
rollback = Pledge
  { ret = (), pre = Single (Atom "beginTx" [])
  , post   = Single (Atom "rollback" [])
  , future = universe }

-- Bad future: transaction never closed
openTx = do { beginTx }

-- Good: transaction properly committed
committedTx = do { beginTx; commit }

-- Good: transaction properly rolled back
abortedTx = do { beginTx; rollback }
\end{lstlisting}

\noindent\lstinline{openTx} retains the full disjunction
$\finally{\mathtt{commit}} \vee \finally{\mathtt{rollback}}$ as its
residual future. \lstinline{committedTx} and \lstinline{abortedTx}
fully discharge the obligation, yielding $\univ$.

\paragraph{RE summary.}
Table~\ref{tab:examples-re} summarises the three RE domains.

\begin{table}[ht]
\caption{RE instance: bad programs and residual annotations.}
\label{tab:examples-re}
\small
\begin{tabular}{llll}
\toprule
Domain & Bad program & Residual \lstinline{future} & Residual \lstinline{pre} \\
\midrule
Memory      & \lstinline{missingFree}           & $\finally{\mathtt{free}(2)}$ & $\univ$ \\
Memory      & \lstinline{freeWithoutMalloc}     & $\univ$ & $\varnothing$ \\
Memory      & \lstinline{doubleFreeImmediate}   & $\varnothing$ & $\univ$ \\
Memory      & \lstinline{leakAndDoubleFree}     & $\varnothing$ & $\univ$ \\
Memory      & \lstinline{allocReturnAddr}       & $\finally{\mathtt{free}(1)}$ & $\univ$ \\
Memory      & \lstinline{allocTwoReturnPair}    & $\finally{\mathtt{free}(1)}\wedge\finally{\mathtt{free}(2)}$ & $\univ$ \\
Mutex       & \lstinline{acquire 1} leaked      & $\finally{\mathtt{release}(1)}$ & $\univ$ \\
Mutex       & release without acquire           & $\univ$ & $\varnothing$ \\
Transaction & no commit/rollback                & $\finally{\mathtt{commit}}\vee\finally{\mathtt{rollback}}$ & $\univ$ \\
\bottomrule
\end{tabular}
\end{table}

\subsection{Bounded Counter (\texorpdfstring{$\GRE$}{GRE} Instance)}
\label{sec:ex-bounded}

The bounded counter (\S\ref{sec:guardedre-instance}) demonstrates that
both the Presburger and RE components are necessary.
The RE alone would accept overflow sequences; the Presburger predicate
alone would accept arbitrarily ordered events.

\begin{lstlisting}
-- Good: init, two increments, one decrement, snapshot
normalUse :: Pledge GuardedRE ()
normalUse = do
  initCounter 0; increment 0 10; increment 0 10
  decrement 0;   snapshot 0

-- Bad: 11 increments against max 10
overflow :: Pledge GuardedRE ()
overflow = do
  initCounter 0
  sequence_ (replicate 11 (increment 0 10))
  snapshot 0
\end{lstlisting}

\noindent
For \lstinline{overflow}, the combined \lstinline{pre} after composing all
eleven increments normalises to
$[\mathbf{false}] \wedge \univ$: the Presburger side collapses to
\lstinline{PFalse} while the RE side remains the universal language,
pinpointing the arithmetic violation.

\paragraph{$\GRE$ summary.}
\begin{center}
\begin{tabular}{lll}
\toprule
Bad program & Residual \lstinline{pre} ($\GRE$) & Cause \\
\midrule
\lstinline{overflow}  & $[\mathbf{false}]\;\univ$ & $h[0] < 10$ violated \\
\lstinline{underflow} & $[\mathbf{false}]\;\univ$ & $h[0] > 0$ violated \\
\lstinline{noInit}    & $[\mathbf{true}]\;\varnothing$   & init event missing \\
\bottomrule
\end{tabular}
\end{center}

\subsection{Task Scheduling (\texorpdfstring{$\WRE$}{WRE} Instance)}
\label{sec:ex-sched}

The task-scheduling example shows that the tropical semiring naturally
accumulates minimum costs.
Each \lstinline{submitJob} operation carries both a submission cost and
a future condition requiring that every submitted job eventually completes.

\begin{lstlisting}
-- Good: submit two jobs, complete both
goodSchedule :: Pledge (WRE Tropical) ()
goodSchedule = do
  submitJob 1 2.0; submitJob 2 3.0
  completeJob 1;   completeJob 2

-- Bad: job 2 never completed
missingComplete :: Pledge (WRE Tropical) ()
missingComplete = do
  submitJob 1 2.0; submitJob 2 3.0
  completeJob 1
\end{lstlisting}

\noindent
For \lstinline{goodSchedule}, \lstinline{wNullable (evalFuture ...)}
returns \lstinline{Tropical 0}, indicating all obligations are
discharged at zero residual cost.
For \lstinline{missingComplete}, the residual future still contains
\lstinline{wFinally (Tropical 0) (complete 2)}, and the overall weight
is $+\infty$, flagging job~2 as unfinished.

The probability-semiring memory example follows analogously: a call
\lstinline{malloc a p} with $p = 0.9$ leaves a residual weight of
$0.9$ if \lstinline{free a} is never called, dropping to $0$ when it is,
matching the reading given in \S\ref{sec:wre-instance}.

\subsection{Heap Memory (SL Instance)}
\label{sec:ex-heap}

The heap-memory example mirrors the RE memory example but uses heap
ownership rather than event traces. Unlike the earlier definitions in this
paper, \lstinline{alloc} and \lstinline{free} here use \lstinline{Top}
futures because, in this particular example, no cell's disposal is a
precondition for allocating another; the obligation of interest is entirely
in \lstinline{post} (what is currently owned), which is what the leak
example below exercises.

\begin{lstlisting}
alloc :: Addr -> Val -> Pledge SL ()
alloc a v = Pledge
  { ret = (), pre = Top, post = Cell a v, future = Top }

free :: Addr -> Val -> Pledge SL ()
free a v = Pledge
  { ret = (), pre = Cell a v, post = Emp, future = Top }

writeCell :: Addr -> Val -> Val -> Pledge SL ()
writeCell a old new = Pledge
  { ret = (), pre = Cell a old, post = Cell a new, future = Top }
\end{lstlisting}

\noindent
\textbf{Good: \lstinline{alloc 0 1 >> alloc 1 2}.}
The combined \lstinline{post} is
$\mathtt{SepStar}\;(\mathtt{Cell}\;0\;1)\;(\mathtt{Cell}\;1\;2)$:
two disjoint ownerships composed by separating conjunction,
correctly reflecting that the two cells reside at distinct addresses.

\noindent
\textbf{Bad: leak (\lstinline{alloc 0 1 >> alloc 1 2 >> free 0 1}).}
After freeing only address~0, the \lstinline{post} retains
$\mathtt{Cell}\;1\;2$, naming the unreleased ownership at address~1. To be
precise about what ``bad'' means for this particular program: the residual
\lstinline{post} of $\mathtt{Cell}\;1\;2$ is not itself a defect (it simply
and correctly records that address~1 is still owned); the point of the
example is diagnostic, showing that \lstinline{Pledge}'s residual, after
composing exactly these three operations, names address~1 and not address~0
as the one still outstanding, which is the piece of information a
leak-detection tool built on top of this instance would want. This is
different in kind from the RE memory examples of
Table~\ref{tab:examples-re}, where a non-\lstinline{universe} residual is
unambiguously a violation of a stated future obligation; here, since
\lstinline{alloc}'s future is \lstinline{Top}, there is no such stated
obligation, and it is the \emph{caller}, inspecting \lstinline{post}, who
decides what counts as a leak.

\noindent
\textbf{Bad: read without ownership (\lstinline{readCell 0 42}).}
The \lstinline{pre} of the bare read is $\mathtt{Cell}\;0\;42$,
which is never discharged, flagging the missing ownership.

\subsection{Bank Account with Presburger Guards (SL Instance)}
\label{sec:ex-bank}

The bank account example demonstrates \lstinline{Pure} Presburger
constraints inside \lstinline{Conj}, and, unlike the heap-memory example
above, gives \lstinline{withdraw} a non-trivial future condition: once
money has been withdrawn against an account, the account must eventually
be reconciled by a \lstinline{closeAccount} or a further balancing
\lstinline{deposit}, so that a withdrawal cannot simply be forgotten about.

\begin{lstlisting}
withdraw :: Addr -> Val -> Val -> Pledge SL ()
withdraw a old amount = Pledge
  { ret = ()
  , pre  = Conj (Pure (PGe (ValAt a) (Lit amount))) (Cell a old)
  , post = Cell a (old - amount)
  , future = \_ -> SepStar (Wand (Cell a (old - amount)) Top)
                           (Pure PTrue) }  -- must eventually reconcile a's cell

closeAccount :: Addr -> Val -> Pledge SL ()
closeAccount a bal = Pledge
  { ret = ()
  , pre  = Conj (Pure (PEq (ValAt a) (Lit 0))) (Cell a bal)
  , post = Emp, future = Top }
\end{lstlisting}

\noindent
The precondition of \lstinline{withdraw} is a \lstinline{Conj} of a pure
Presburger guard ($h[a] \geq \mathtt{amount}$) and a spatial ownership
claim ($\mathtt{Cell}\;a\;\mathtt{old}$).
Without the \lstinline{Pure} constructor, \lstinline{Conj} of two purely
spatial predicates would be either trivial or unsatisfiable (see
\S\ref{sec:sl-instance}); the combination of \lstinline{Pure} and
\lstinline{Conj} is what makes value-dependent preconditions expressible.

\noindent
\textbf{Bad: overdraft.}
After \lstinline{deposit 0 0 100}, the account holds 100.
The call \lstinline{withdraw 0 100 150} carries
$\pre = \mathtt{Conj}\;(\mathtt{Pure}\;(h[0] \geq 150))\;(\mathtt{Cell}\;0\;100)$.
Since $100 < 150$, the pure guard is not satisfiable, and the residual
\lstinline{pre} of the composed program retains this unsatisfiable constraint.

\noindent
\textbf{Bad: unreconciled withdrawal.}
\lstinline{deposit 0 0 100 >> withdraw 0 100 40} leaves the residual future
as $\mathtt{SepStar}\;(\mathtt{Wand}\;(\mathtt{Cell}\;0\;60)\;\top)\;(\mathtt{Pure}\;\mathbf{true})$,
recorded symbolically per \S\ref{sec:sl-instance}: composing a subsequent
\lstinline{closeAccount 0 60} discharges it, matching the RE instance's
disjunctive-obligation pattern from the transaction example
(\S\ref{sec:examples}), but here over heap ownership instead of a trace.

\subsection{Linked List (SL Instance)}
\label{sec:ex-list}

A singly-linked-list node occupies two consecutive cells:
address $a$ holds the payload value and address $a{+}1$ holds the next pointer.

\begin{lstlisting}
allocNode :: Addr -> Val -> Val -> Pledge SL ()
allocNode a val next = Pledge
  { ret = (), pre = Top
  , post = SepStar (Cell a val) (Cell (a+1) next)
  , future = Top }
\end{lstlisting}

\noindent
\textbf{Two-node list (\lstinline{allocNode 0 10 2 >> allocNode 2 20 (-1)}).}
The combined \lstinline{post} is a four-cell \lstinline{SepStar}:
\[
  (\mathtt{Cell}\;0\;10 \sep \mathtt{Cell}\;1\;2)
  \sep
  (\mathtt{Cell}\;2\;20 \sep \mathtt{Cell}\;3\;{-1})
\]
This is the canonical representation of a two-node list in separation logic:
all four cells are disjointly owned, and the structure of \lstinline{SepStar}
mirrors the structure of the list.

\paragraph{SL summary.}
Table~\ref{tab:examples-sl} summarises the four SL domains.
The key difference from the RE instance is that residual annotations are
heap predicates rather than traces: a leaked cell appears in \lstinline{post},
an unsatisfiable ownership constraint in \lstinline{pre}, and, for the bank
account, an outstanding wand in \lstinline{future}.

\begin{table}[ht]
\caption{SL instance: bad programs and residual annotations.}
\label{tab:examples-sl}
\begin{tabular}{lllll}
\toprule
Domain & Bad program & Residual \lstinline{post} & Residual \lstinline{pre} & Residual \lstinline{future} \\
\midrule
Heap memory  & \lstinline{Cell 1 2} not freed  & $\mathtt{Cell}\;1\;2$ (owned) & $\top$ & $\top$ \\
Heap memory  & read without ownership           & $\mathbf{emp}$ & $\mathtt{Cell}\;0\;42$ & $\top$ \\
Bank account & insufficient balance             & $\mathtt{Cell}\;0\;100$ & unsatisfiable guard & $\top$ \\
Bank account & withdrawal not reconciled        & $\mathtt{Cell}\;0\;60$ & $\top$ & $\mathtt{Wand}\;(\mathtt{Cell}\;0\;60)\;\top$ \\
Linked list  & read without ownership           & $\mathbf{emp}$ & $\mathtt{Cell}\;0\;99$ & $\top$ \\
\bottomrule
\end{tabular}
\end{table}

%% =========================================================================
\section{Monad Laws}
\label{sec:monad-laws}
%% =========================================================================

Every monad must satisfy three equational laws that together guarantee
sequential composition behaves predictably:
\begin{itemize}
  \item \textbf{Left identity}: \lstinline{pure x >>= f = f x}.
        Wrapping a value in \lstinline{pure} and immediately binding it is
        the same as applying \lstinline{f} directly; \lstinline{pure}
        introduces no effect of its own.
  \item \textbf{Right identity}: \lstinline{e >>= pure = e}.
        Binding a computation to \lstinline{pure} is the same as the
        computation itself; \lstinline{pure} consumes no observable state.
  \item \textbf{Associativity}:
        \lstinline|(e >>= f) >>= g = e >>= (\x -> f x >>= g)|.
        Re-parenthesising a chain of binds does not change its meaning;
        sequential composition is order-preserving but grouping-independent.
\end{itemize}
When these laws hold, programmers can freely refactor monadic chains
without changing the observable semantics: extracting sub-expressions,
fusing steps, or reordering \lstinline{do}-blocks are all valid.

\paragraph{Scope of the laws: the $(\ret,\post,\fut)$ triple, not \texttt{pre}.}
We state upfront, rather than only at the end of this section, exactly what
is and is not proved. \lstinline{Pledge}'s four fields do not \emph{all}
satisfy the monad laws under ordinary equality: \lstinline{pre} follows the
Hoare-rule discipline of Equation~\eqref{eq:pre-bind}, which deliberately
deviates from left-identity when a precondition violation is detected (a
violation must remain visible after further binding, not be silently
absorbed). What we prove below is that the triple $(\ret, \post, \fut)$,
i.e.\ \lstinline{Pledge} values compared \emph{ignoring} the \lstinline{pre}
field, satisfies the three laws exactly. Equivalently: \lstinline{Pledge} is
a lawful \lstinline{Monad} once \lstinline{Eq} (or the informal notion of
``the same specification'') is defined to ignore \lstinline{pre}, and
\lstinline{pre} is carried alongside it as a separate, non-monadic
contract-checking annotation, analogous to a writer-monad log that is
inspected but not required to respect the monad laws on its own. We take
this position, rather than trying to force \lstinline{pre} into the same
equality, because absorbing precondition violations under left-identity
would make \lstinline{Pledge} silently forget a detected contract violation
whenever it is immediately rebound, which is the opposite of what a
specification-checking tool should do.

Because \lstinline{future} has type \lstinline{a -> eff}, the laws are
stated in terms of \lstinline{evalFuture e = future e (ret e)}, written
$\widehat{\fut}(e)$ for brevity.
The laws hold for the triple $(\ret, \post, \widehat{\fut})$ assuming
\lstinline{Composable} forms a monoid under \lstinline{<>} with unit
\lstinline{empty}, and \lstinline{universe} is the unit for \lstinline{(/\\)}.

\paragraph{Left identity.} \lstinline{pure x >>= f = f x}.

$\ret$: $\ret(f\,x)$. \checkmark \quad
$\post$: $\varepsilon \mathbin{<>} \post(f\,x) = \post(f\,x)$. \checkmark

For the future, let $fe = f(\ret(\mathtt{pure}\;x)) = f\,x$.
Since $\widehat{\fut}(\mathtt{pure}\;x) = \univ$, subtraction's
$\Sigma^*$-base-case gives $\post(fe) \setminus \univ = \varepsilon$
directly from the \lstinline{Composable} instance's second base clause
(\S\ref{sec:re-instance}); the identity used here is exactly
$r_1 \setminus \univ = \varepsilon$ read with $r_1$ replaced by $\post(fe)$,
i.e.\ \emph{any} postcondition, once compared against the universal
obligation, discharges completely, because there is nothing left to name.
Hence:

\begin{align*}
  \widehat{\fut}(\mathtt{pure}\;x \mathbin{>\!\!>\!\!=} f)
    &= \bigl(\post(fe) \setminus \widehat{\fut}(\mathtt{pure}\;x)\bigr)
       \wedge \widehat{\fut}(fe) \\
    &= \bigl(\post(f\,x) \setminus \univ\bigr)
       \wedge \widehat{\fut}(f\,x)
     \;=\; \varepsilon \wedge \widehat{\fut}(f\,x)
     \;=\; \widehat{\fut}(f\,x). \;\checkmark
\end{align*}

\paragraph{Right identity.} \lstinline{e >>= pure = e}.

$\ret$: $\ret(e)$. \checkmark \quad
$\post$: $\post(e) \mathbin{<>} \varepsilon = \post(e)$. \checkmark

Let $fe = \mathtt{pure}(\ret(e))$, so $\post(fe) = \varepsilon$ and
$\widehat{\fut}(fe) = \univ$.
Then, using the other base case $\varepsilon \setminus r_2 = r_2$:

\begin{align*}
  \widehat{\fut}(e \mathbin{>\!\!>\!\!=} \mathtt{pure})
    &= \bigl(\varepsilon \setminus \widehat{\fut}(e)\bigr)
       \wedge \univ
     \;=\; \widehat{\fut}(e) \wedge \univ
     \;=\; \widehat{\fut}(e). \;\checkmark
\end{align*}

\paragraph{Associativity.}
Re-parenthesising \lstinline{(e >>= f) >>= g} as
\lstinline|e >>= (\x -> f x >>= g)| must not change the specification.

$\ret$ and $\post$: both sides yield $\ret(g\,r_2)$ and
$\post(e) \mathbin{<>} \post(f\,r_1) \mathbin{<>} \post(g\,r_2)$,
following from associativity of \lstinline{<>}. \checkmark

For the future, let $r_1 = \ret(e)$, $r_2 = \ret(f\,r_1)$.
Unfolding the bind formula (Equation~\eqref{eq:future-bind}) on both sides
and applying associativity of $\wedge$, both reduce to the same expression:
\[
  \bigl(\post(g\,r_2) \setminus
        \bigl(\post(f\,r_1) \setminus \widehat{\fut}(e)\bigr)\bigr)
  \;\wedge\; \widehat{\fut}(f\,r_1)
  \;\wedge\; \widehat{\fut}(g\,r_2). \;\checkmark
\]

\paragraph{Precondition and the monad-law status of \lstinline{pre}, revisited.}
Returning to the point flagged at the start of this section: the residual
$\post(e) \setminus \pre(e_2)$ of Equation~\eqref{eq:pre-bind} captures
which part of $e_2$'s precondition $e$'s postcondition has not discharged;
full coverage gives the standard Hoare sequencing rule, and absent coverage
collapses to $\varnothing$, flagging the violation immediately and
\emph{persistently} through further binds, which is the behaviour we want
from a contract checker and which is incompatible with unrestricted
left-identity on \lstinline{pre}. In practice, a programmer inspects
\lstinline{pre} and \lstinline{evalFuture} after constructing the
specification term, and both fields equal \lstinline{universe} if and only
if the composed program is contract-correct (Theorem~\ref{thm:soundness}
for the $\RE$ instance).

%% =========================================================================
\section{Related Work}
\label{sec:related}

\paragraph{Future conditions and temporal verification.}
The proposal this paper implements is due to \citet{song2025futurecond},
who introduce future conditions as a first-class extension to pre/post
contracts and formalise the notion that an operation may discharge temporal
obligations onto its continuation, at the level of a formal system rather
than a runnable library. \citet{song2024provenfix} apply a related
regular-expression encoding of temporal properties in ProveNFix, to guide
automated program repair rather than to specify obligations directly.
Pledge's relationship to both is one of realisation and extension: we take
the future-condition idea of \citet{song2025futurecond} and give it a
concrete monadic interface in Haskell, a derivative-based
\lstinline{subtraction} operator that implements their discharge relation
algorithmically (\S\ref{sec:deriv}), a proof that this operator is sound
with respect to trace membership (\S\ref{sec:soundness}), and four
\lstinline{Composable} instances --- Presburger-guarded traces, semiring
weights, and separation-logic heaps, alongside plain traces --- that
neither prior paper instantiates. Where \citet{song2024provenfix} use
temporal specifications to drive automated repair, Pledge is oriented
toward identifying violations through annotation inspection, prior to any
effectful execution, leaving repair to the programmer or to a separate tool.

\paragraph{Pre/post-condition contracts in Haskell.}
\emph{Liquid Haskell}~\cite{vazou2014} expresses pre/post-conditions as
refinement types verified by an SMT solver.
Findler and Felleisen~\cite{findler2002} establish a runtime contract
discipline for higher-order functions.
QuickCheck~\cite{claessen2000} supports conditional properties via the
implication combinator \emph{==>}, which restricts a property to
inputs satisfying a given precondition.
Swierstra and Baanen~\cite{swierstra2019} embed Hoare triples as an indexed
monad, propagating pre/post annotations through monadic bind analogously to
\lstinline{Pledge}.
Pledge complements these by adding a \emph{temporal} dimension through the
future condition, without requiring type-system extensions.

\paragraph{Parameterised and graded monads.}
\citet{atkey2009} introduces parameterised monads, written
$T(S_1, S_2, A)$, for Hoare-logic-style pre/post state tracking.
\citet{katsumata2014} generalises this to graded monads over a
monoid; \citet{orchard2016} show effect systems and session
types are instances.
Pledge is related but orthogonal: rather than indexing by effect \emph{type}
at the type level, we carry an explicit obligation as
\lstinline{future :: a -> eff} at the value level, preserving compatibility
with the standard \lstinline{Monad} class and \lstinline{do}-notation, in
the specific, narrower sense of \S\ref{sec:effects}: composing
specifications, not composing specifications together with the real effects
they describe.

\paragraph{Resource management in Haskell.}
The \lstinline{bracket} combinator lexically scopes cleanup but cannot
propagate obligations beyond the lexical scope.
The \lstinline{ResourceT}~\cite{snoyman2011resourcet} library threads a
finaliser queue through every bind, discharging cleanup when the block
exits.
Both \lstinline{Acquire} and \lstinline{Managed} generalise this to an
applicative/monadic interface.
All three handle the \emph{cleanup} half of a resource obligation but
cannot express \emph{protocol} or \emph{quantitative} obligations,
which require explicit future-condition propagation across bind steps.

\paragraph{Linear types.}
Linear Haskell~\cite{bernardy2018linear} introduces multiplicity annotations
enforced by GHC since version 9.0; \lstinline{linear-base}~\cite{linearbase2021}
provides a \lstinline{RIO} monad with linear file handles, statically
guaranteeing every opened handle is consumed exactly once.
Linear types enforce \emph{structural} obligations (use exactly once),
whereas future conditions enforce \emph{temporal} obligations (use in a
specified order, or with a specified probability).
A hybrid combining linear types for uniqueness with Pledge for temporal
ordering could achieve stronger guarantees than either alone.

\paragraph{Effect systems and algebraic effect handlers.}
Algebraic effect handlers~\cite{pretnar2015} decompose effectful computations
into an operation set and a separate handler; resource lifecycle obligations
are typically encoded in the handler's state machine.
Polysemy's \lstinline{Scoped} effect attaches resource creation and cleanup
to a lexical scope threaded through the effect row.
Pledge takes the complementary route described in \S\ref{sec:effects}:
obligations are values in the annotation fields of a pure specification
term, checked as a shadow computation, rather than a handler intercepting
a live effect row.

\paragraph{Session types.}
Session types~\cite{honda1993,gay2010} enforce communication protocols
at the type level.
The \lstinline{sessiontypes} library~\cite{sessiontypes2017} realises this
via an indexed monad \lstinline{STTerm m s s' a}; the type checker rejects
programs that violate the protocol order.
Pledge is more lightweight: it operates within the standard
\lstinline{Monad} class, admits quantitative and heap-ownership obligations
that session types do not, and is applicable beyond message-passing.

\paragraph{Typestate.}
Typestate systems~\cite{strom1986,bierhoff2007} encode resource lifecycle
protocols in the type system.
Brady's \lstinline{ST} monad encoding in Idris~\cite{brady2013idris}
demonstrates full compile-time enforcement in a dependently typed language.
Pledge achieves comparable expressiveness at the library level using
standard Haskell type classes, at the cost of moving checking to the point
where the specification term's residual annotations are inspected, rather
than to type-inference time: violations surface as non-trivial residual
annotations in \lstinline{pre} or \lstinline{future}, which the programmer
evaluates before executing effects rather than relying on the type checker
to reject programs.

\paragraph{Runtime monitoring.}
Runtime verification~\cite{leucker2009} checks temporal properties by
instrumenting a concrete execution.
Pledge propagates obligations in the specification layer: residuals
surface as soon as annotations are composed, prior to any effectful run,
without requiring a concrete execution, at the cost --- made explicit in
\S\ref{sec:effects} and \S\ref{sec:discussion} --- of not itself catching
divergence between the specification and the real code it shadows.

\paragraph{Brzozowski derivatives.}
Derivatives of regular expressions~\cite{brzozowski1964} have found
application in parsing~\cite{owens2009,might2011} and model checking.
Equation~\eqref{eq:deriv-not} extends this to complement via
$\partial_a(\neg r) = \neg(\partial_a r)$; our subtraction operator is a
novel application to future-condition propagation in a monadic setting,
made precise by the soundness theorem of \S\ref{sec:soundness}.

%% =========================================================================
\section{Discussion and Future Work}
\label{sec:discussion}
%% =========================================================================

\paragraph{Decidability.}
Equality of normalised regular expressions is decidable, so the
fixed-point computation in \lstinline{subtraction} terminates whenever
the derivative sequence is finite, which holds for the regular
languages used in our examples. Extending to context-free or
$\omega$-regular specifications remains future work.

\paragraph{Precondition monad laws.}
As discussed at the start of \S\ref{sec:monad-laws}, the \lstinline{pre}
field satisfies the Hoare-rule property rather than strict monad-law
equality. A natural refinement is to make \lstinline{pre} a separate,
non-monadic annotation computed by a static analysis pass over the
\lstinline{Pledge} term, separating the behavioural monad from the
precondition contract checker; we have not implemented this refactoring,
and its main open question is whether the resulting two-pass design still
produces the same immediate, per-bind violation signal that the current
design gives.

\paragraph{Closing the gap between the specification and real effects.}
\S\ref{sec:effects} describes the shadow-specification pattern as the
paper's intended use, and is explicit that \lstinline{Pledge} itself does
not check that the shadow specification and the real \lstinline{IO}
computation stay in sync, nor does it intervene if an asynchronous
exception causes the real code to diverge from its specification mid-run.
Lifting Pledge into an effect-handler framework (\texttt{effectful} or
\texttt{polysemy}), so that a \lstinline{Pledge} annotation is attached
directly to each operation's real effect handler and updated automatically
as the effect executes, would close this gap and remove the possibility of
drift between the two. A lighter-weight intermediate step, requiring no new
effect-handler integration, is to pair \lstinline{Pledge} specifications
with \lstinline{bracket} or \lstinline{mask} at the call sites already
identified by the specification's primitives, so that at least the
obligations \lstinline{Pledge} names are also runtime-enforced at those
points; we regard this as the most promising near-term direction and leave
it for future work.

\paragraph{Closed-world alphabet assumption.}
The \lstinline{firstWith} function uses
$\Sigma_0 = \mathit{atoms}(r_1)\cup\mathit{atoms}(r_2)$ as its working
alphabet.
This is sound under the \emph{closed-world assumption} that every relevant event
is mentioned in the two operands.
A fully open-world implementation would accept an explicit $\Sigma$ declaration,
enabling complement over an unrestricted domain.

\paragraph{Connection to indexed and graded monads.}
Because $\mathtt{future} :: a \to \mathit{eff}$ indexes the obligation by the
return value, Pledge echoes the indexed/graded-monad literature where types are
indexed by computational effects; here the index is a return value instead of
an effect row.

\paragraph{Toward solver-backed SL discharge.}
\S\ref{sec:sl-instance} is explicit that the SL instance's
\lstinline{subtraction} records the magic wand symbolically, without
consulting a heap model or invoking a solver, and \S\ref{sec:soundness}
notes that Theorem~\ref{thm:soundness}'s style of guarantee therefore does
not currently extend to it. Wiring the residual \lstinline{SL} predicates
into an existing decision procedure (Z3 via \lstinline{sbv}, or a dedicated
separation-logic prover such as Smallfoot) so that \lstinline{Wand} nodes
are checked for satisfiability rather than merely recorded is the concrete,
scoped next step toward closing this gap, and toward a soundness statement
for $\SL$ analogous to Theorem~\ref{thm:soundness}.

\paragraph{Parallel composition.}
The \lstinline{Pledge} monad is inherently sequential; concurrent programs
require a parallel combinator $e_1 \parallel e_2$.
In the RE instance this corresponds to the shuffle product $r_1 \shuffle r_2$;
the SL instance admits a parallel connective via the concurrent separation
logic frame rule~\cite{o2007resources}.
Extending \lstinline{Composable} with a \lstinline{parallel} operation is
a natural next step.

\paragraph{Quantitative future conditions.}
A natural further direction is to connect $\WRE$ to probabilistic temporal
logics (PCTL) or quantitative Kleene algebras; in the SL instance, a
quantitative analogue could express amortised resource bounds or expected heap
usage.

%% =========================================================================
\section{Conclusion}
\label{sec:conclusion}
%% =========================================================================

We have presented \emph{Pledge}, a Haskell library that realises
\citet{song2025futurecond}'s proposal of \emph{future conditions} as a
concrete monad: annotations that constrain what the remainder of a program
must accomplish, checked as a pure specification before, or independently
of, any real effectful execution.

The central insight is that trace-based, arithmetic, quantitative, and
heap-ownership obligations are structurally identical at the level of a
five-operation \lstinline{Composable} algebra.
All four instances ($\RE$, $\GRE$, $\WRE$, $\SL$) share a single monadic
bind formula, and the diagnostic signal is uniform: a non-trivial residual
in \lstinline{pre} or \lstinline{future} identifies exactly which obligation
remains unmet and why, a diagnostic we prove sound for the $\RE$ instance
(Theorem~\ref{thm:soundness}).

Pledge requires no language extensions beyond standard type classes and
composes via a standard \lstinline{Monad} instance, letting a specification
term be built with ordinary \lstinline{do}-notation; as discussed in
\S\ref{sec:effects}, this composition operates on the pure specification
layer itself; connecting that layer to a real, running effectful
computation is the shadow-specification pattern we describe, and tightening
that connection is the paper's main direction for future work, alongside
parallel composition (shuffle product in the RE instance; frame rule in the
SL instance) and solver-backed residual discharge for the SL instance.

%% =========================================================================
\bibliographystyle{ACM-Reference-Format}
\bibliography{references}

\end{document}